% !TEX root = ../main.tex
\chapter{연구 방법}
\label{ch:methodology}

\ref{ch:literature}장은 다섯 가지 빈틈(아직 채워지지 않은 연구 부분)으로 끝났어요. 이 장에서는 그 빈틈마다 잴 수 있는 절차를 하나씩 정해요. 또 엔도스랭크(EndorseRank), AWP(적응형 가중 페이지랭크), 그리고 이 둘의 두 층(한 종류 화살표로 만든 그래프)으로 만든 결합 걷기(화살표를 따라 점에서 점으로 옮겨 다니는 계산)들을 아비트럼 원(Arbitrum One)에서 평판 점수로 어떻게 평가하는지 밝혀요. 엔도스랭크는 ERC-20(토큰 공통 규칙) 허락(allowance, 토큰을 얼마까지 대신 써도 된다는 허락) 그래프(점과 화살표로 이은 그림) 위에서 돌아가고, AWP는 송금(transfer) 그래프 위에서 돌아가요\cend{3}. 둘 다 AWP를 만든 사람들이 정한 대로 화살표(엣지)에 무게(가중치)를 달고, 무게를 단 페이지랭크(PageRank)로 점수를 구해요\cend{15}. AWP는 송금을 보내는 활동에 비례해서(많을수록 그만큼 더) 걷기를 다시 출발(재시작)시키고, 엔도스랭크는 모든 점(노드)에서 고르게 다시 출발시켜요.

방법 면에서 이 연구는 구성개념(점수가 재려고 하는 생각) 일치도, 곧 분야 일치도를 보는 방식을 따라요. 평판 점수는 끊김 없이 이어지는 값으로 매긴 순위예요. 그리고 점검은 두 값만 있는 부도 라벨(빌린 돈을 못 갚았는지 아닌지만 적은 정답 값) 대신, 지갑마다 볼 수 있는 통계에 기대요\cend{22}. 순위를 매기는 방법 하나하나는, 직접 볼 수 없는 평판이라는 구성개념을 보여 주는 지표 하나예요. 그래프 일관성(점수가 자기 그래프의 통계와 잘 맞는지)은 같은 기간의 송금 통계와 허락 통계로 점검해요. 기간 밖 점검은 시간 홀드아웃(나중 기간을 떼어 두고 맞혀 보는 시험)이에요. 여기서는 앞으로 주인(토큰 주인)과 스펜더(spender, 허락을 받아 대신 쓰는 쪽) 사이에 새로 생기는 허락, 그리고 새 송금자(새로 송금을 보내온 주소)를 라벨로 써요. 이 연구의 주장은 비교하는 주장이에요. 모든 점수를 각 구성개념의 점검과 견주고, 그 점수가 매끄럽게 다듬는(평활하는) 원시 차수(화살표를 그냥 센 수)와도 견주어요. 거래 이익이나 신용 부도를 예측하는(미리 맞히는) 것은 아니에요.

모든 절차는 정해 둔 관찰 기간 안의, 누구나 볼 수 있는 블록체인 로그(계약이 남기는 기록)로 해요. 시드(무작위 계산의 시작 값)는 고정되어 있고, 평가 처리 과정(파이프라인)은 공개되어 있어요. 그래서 자료 꺼내기(추출)를 다시 하는 사람이나, 저장해 둔 parquet 파일을 가진 사람은 누구나 이 논문에 실린 표를 다시 만들 수 있어요.

실제 작업에서 무엇을 고를 때마다 네 가지 원칙을 따라요.
\begin{itemize}
\item 같은 씨앗 지갑으로 비교하기: 모든 점수를 똑같은 지갑 모음에서 계산해요. 그래서 점수 사이의 어떤 차이도 표본 뽑기 때문에 생길 수 없어요.
\item 미리 정해 두기: 관찰 기간, 토픽 해시(이벤트 종류를 알려 주는 고유한 값), 감쇠 계수(화살표를 따라갈 확률), 시간 감쇠의 매개변수(미리 정하는 값), 부트스트랩(다시 뽑기) 횟수와 시드, 벤치마크(속도 재기) 시드를 버전을 매겨 관리하는 설정에 정해 두었어요. 결합 연산자(두 층을 함께 걷게 하는 계산 규칙)와 그 무게들, 그리고 RQ5(다섯 번째 연구 질문)에 대한 첫 번째 규칙은 2026년 9월 14일에 더했어요. 결합 점수를 하나라도 계산하기 전이었어요. RQ5에 대한 두 번째 규칙과 그 규칙의 6월부터 8월까지 기간은 2026년 10월 1일에 등록(계획을 미리 공개해 두기)했어요. 그 기간의 로그를 하나라도 꺼내기 전이었어요(\ref{sub:fresh-holdout}절). 엔도스랭크와 AWP는 같은 기간 점검 결과와 봄 라벨을 이미 안 뒤에 점수를 매겼어요. 그래서 C-PR(결합 페이지랭크, 두 그래프를 함께 걷는 점수)을 AWP와 견주는, 등록해 둔 민감도 분석(조건을 바꿔 결과가 얼마나 움직이는지 보기) 하나를 빼면, 이 둘이 들어간 비교는 어느 것도 미리 정해 두지 않았어요. 모든 비교는 신뢰구간(믿을 만한 범위)과 함께 보고해요.
\item 만들기와 시간 재기: 코호트(묶음)마다 화살표 목록은 어떤 시간 재기도 시작하기 전에 한 번만 만들어요. 벤치마크는 풀이 프로그램(solver)을 부르는 데 걸린 시간을 재요.
\item 범위: 가설(H) H1과 주장(C) C3은 엔도스랭크와 AWP가 매긴 순서를 비교해요. 나머지 가설들과 주장 C1, C2, C4는 모든 점수에 대한 것이고, C5는 결합 연산자와 그 빼 보기 실험(ablation, 한 부분을 빼거나 바꿔 보는 실험)에 대한 것이에요. RQ4와 RQ5는 스펜더 홀드아웃 코호트를 쓰고, 등록 재현(계획을 먼저 공개해 두고 새 자료로 다시 해 보기)에는 트레이더(거래하는 사람) 코호트가 더해져요.
\end{itemize}

\dissertationSubheading[sec:notation]{기호}

이 장과 \ref{ch:results}장, \ref{ch:discussion}장에서 쓰는 기호는 아래 표~\ref{tab:notation-methods}에 모아 두었어요. 줄인 목록은 표~\ref{tab:notation}에 있어요. 벡터(수 여러 개를 한 줄로 늘어놓은 것)는 굵은 글씨로, 집합(모음)은 멋 글씨체(필기체)로 써요. 지갑 주소는 속을 들여다볼 수 없는 이름표로만 다뤄요. 우리는 주소를 바깥의 어떤 신원(실제로 누구인지)과도 이어 붙이지 않아요.

\begin{table}[htbp]
\centering
\caption{그래프, 점수, 대리 지표(대신 재는 값), 매개변수에 쓰는 기호.}
\label{tab:notation-methods}
\fitwidth{%
\begin{tabular}{lp{0.72\linewidth}}
\toprule
기호 & 뜻 \\
\midrule
$\mathcal{W}$ & 매칭 지갑 코호트(모든 점수가 똑같이 점수를 매기는 지갑 묶음) ($|\mathcal{W}|=n=5{,}521$) \\
$\mathcal{W}_{\mathrm{pool}}$ & 확장 지갑 풀(더 큰 주소 모음) ($|\mathcal{W}_{\mathrm{pool}}|=N=99{,}080$) \\
$G=(V,E)$ & 페이지랭크 풀이 프로그램에 넘기는 그래프: $V$는 남긴 화살표 $E$의 끝점들을 담음 \\
$G_{\mathrm{approve}}=(V_A,E_A,w_A)$ & EndorseRank 허락 그래프 \\
$G_{\mathrm{transfer}}=(V_T,E_T,w_T)$ & AWP 송금 그래프 \\
$a^{\star}(u,v,\mathrm{token})$ & 토큰 하나에서 주인 $u$가 스펜더 $v$에게 기간 안에 준 마지막 허락(원래 단위, 소수점 조정을 하지 않은 가장 작은 단위) \\
$\Delta t^{\star}(u,v,\mathrm{token})$ & $a^{\star}$를 정한 허락의 나이(지난 시간) \\
$w_A(u,v)$ & EndorseRank 화살표 무게, 주인 $u\to$ 스펜더 $v$: 토큰마다 $\sigma(\Delta t^{\star})V(a^{\star})$를 더한 값 \\
$w_T(u,v)$ & AWP 화살표 무게, 보낸 쪽(송신자) $u\to$ 받는 쪽(수신자) $v$: 송금마다 $\sigma(\Delta t_e)V(x_e)$를 더한 값 \\
$x_e$ & 송금 $e$의 토큰 양(원래 단위) \\
$V(z)$, $b$ & AWP의 상한이 있는 값 변환(금액이 아무리 커도 1을 넘지 않게 바꾸는 식), $V(z)=2/(1+e^{-bz})-1$, 원래 단위에서 $b=1$ \\
$X_u$ & 주소 $u$의 활발함(활동 정도), AWP의 재시작 무게 \\
$\tilde w_A(u,v)$, $\tilde w_T(u,v)$ & C-PR의 두 층의 무게, 금액 그대로 \\
$o(u)$; $o_A(u)$, $o_T(u)$ & 점 $u$의 나가는 무게의 합 $\sum_v w(u,v)$; 허락 층과 송금 층에서의 같은 값 \\
$P$ & 이동 행렬(전이 행렬), $P_{u,v}=w(u,v)/o(u)$; 매달린 점(나가는 화살표가 없는 점)의 행은 모두 0 \\
$\mathbf{1}[\cdot]$ & 지시 함수: 조건이 맞으면 $1$, 아니면 $0$ \\
$\mathbf{1}_{\mathrm{dang}}$ & 매달린 점을 표시하는 지시 벡터, 각 칸은 $\mathbf{1}[o(u)=0]$ \\
$\mathbf{r}$, $\mathbf{r}^{(t)}$ & 페이지랭크 벡터와, $t$걸음 반복한 뒤의 그 값 \\
$\boldsymbol{\pi}$ & 재시작(순간 이동) 분포(어디서 다시 출발하는지의 비율), 매달린 점의 몫(질량, 그 점에 머무는 확률)에도 씀: EndorseRank와 C-PR은 고르게, AWP는 $X_u$에 비례, S-PR(보증 씨앗 페이지랭크)은 허락 층의 점수를 $V_T$로 좁혀 합이 1이 되게 다시 맞춘 것 \\
$d$ & 페이지랭크 감쇠 계수(기본값 $0.85$) \\
$\varepsilon$ & 페이지랭크 수렴(값이 더는 거의 바뀌지 않게 됨) 허용 오차($10^{-8}$) \\
$I_{\max}$ & 페이지랭크 반복(한 바퀴)의 최대 횟수($300$) \\
$\lambda$ & 결합 연산자 C-PR에서 허락 층의 무게, 곧 섞는 비율($\lambda=0.5$ 주 분석; $0.25$, $0.75$ 민감도 분석; $\lambda=1$과 $\lambda=0$에서는 한 층만 걸음) \\
$\alpha_A(u)$, $\alpha_T(u)$ & 점 $u$의 C-PR 행에서 각 층의 무게; $D_{\lambda}(u)$는 둘의 공통 분모 \\
$P^{(\lambda)}$ & C-PR의 결합 이동 행렬 \\
$\mathbf{s}_M$ & 방법 $M$의 평판 점수 벡터 \\
$\mathbf{p}_q$ & 대리 지표 $q$의 값 벡터 \\
$\sigma(\Delta t)$ & 허락이나 송금의 나이에 거는 로지스틱 시간 감쇠(오래된 일일수록 가볍게 세는 무게) \\
$k$, $t_0$ & 감쇠가 가파른 정도($0.01$)와 가운데 점($180$일) \\
$\Delta t$ & 감쇠 기준 시점($T_{\mathrm{obs}}$; 홀드아웃에서는 $t_1$)보다 며칠 전에 일어난 이벤트인지(이벤트의 나이) \\
$T_{\mathrm{obs}}$ & 관찰 끝 시각(2026-05-31 23:59:59 UTC) \\
$t_1$, $t_2$ & 홀드아웃 점수 동결 시각(점수를 고정하는 때, 2026-02-28 23:59:59 UTC)과 라벨 기간의 끝(2026-05-31 23:59:59 UTC) \\
$\rho$ & 스피어먼의 순위 상관(스피어먼의 로, 두 순위가 얼마나 비슷한지 나타내는 수) \\
$R_i$, $S_i$ & 지갑 $i$의 점수와 대리 지표 값의 중간 순위(동점인 값들은 그 자리들의 평균 순위를 함께 가짐) \\
$\tau$, $\tau_b$ & 켄달의 순위 상관(켄달의 타우, 두 줄 세우기가 얼마나 같은 순서인지 나타내는 수); 보고한 값은 모두 동점을 바로잡은 $\tau_b$ \\
$n_c$, $n_d$ & 순서가 맞는 지갑 쌍과 어긋나는 지갑 쌍의 수 \\
$n_0$; $n_1$, $n_2$ & 모든 쌍의 수 $n(n-1)/2$; 첫 번째 순위와 두 번째 순위에서 각각 동점인 쌍의 수 \\
\bottomrule
\end{tabular}
}
\end{table}

따로 말하지 않으면, 페이지랭크는 Page 등\cend{15}이 만든 페이지랭크에 무게를 단 변형을 뜻해요. 반복 $t$에서 값을 새로 고치는 식은 다음과 같아요.
\begin{equation}
\label{eq:pagerank}
\mathbf{r}^{(t+1)}
=
(1-d)\,\boldsymbol{\pi}
+
d\,P^{\top}\mathbf{r}^{(t)}
+
d\,\bigl(\mathbf{r}^{(t)\top}\mathbf{1}_{\mathrm{dang}}\bigr)\,\boldsymbol{\pi},
\end{equation}
여기서 $P$는 아래 식~\eqref{eq:transition-impl}의 이동 행렬이에요. $\boldsymbol{\pi}$는 점들 위의 재시작 분포예요. 방법이 다른 분포를 정하지 않으면 고른 분포($\pi_u=1/|V|$)예요(다른 분포를 쓰는 방법은 \ref{sub:awp-graph}절의 AWP와 \ref{sub:coupled-graph}절의 S-PR이에요). $\mathbf{1}_{\mathrm{dang}}$은 매달린 점들을 표시하는 지시 벡터이고, 각 칸은 $\mathbf{1}[o(u)=0]$이에요. 그래서 세 번째 항은 나가는 화살표가 없는 점들이 가진 몫을 모아요. $w(u,v)$를 $u$에서 $v$로 가는 방향이 있는 화살표의 무게라고 하고, $o(u)=\sum_v w(u,v)$를 나가는 무게의 합이라고 하면, 이동 행렬은 다음과 같아요.
\begin{equation}
\label{eq:transition-impl}
P_{u,v}
=
\begin{cases}
w(u,v)/o(u) & \text{만약 }o(u)>0,\\
0 & \text{그 밖에는.}
\end{cases}
\end{equation}
매달린 점의 행은 모두 0이에요. 그래서 $P$는 준확률 행렬(행의 합이 1보다 작을 수도 있는 행렬)이에요. 그 점들에서 잃은 몫은 식~\eqref{eq:pagerank}의 세 번째 항을 거쳐 돌아오고, $\boldsymbol{\pi}$에 따라 나누어 퍼져요. 0인 행마다 $1/|V|$를 채우면 식~\eqref{eq:transition}의 행 확률 행렬(행마다 합이 1인 행렬)이 돼요. 이때 $\boldsymbol{\pi}$가 고르면, 식~\eqref{eq:pagerank}의 꼴은 식~\eqref{eq:google-matrix}의 구글 행렬 꼴과 같아져요. 반복은 $\mathbf{r}^{(0)}=\mathbf{1}/|V|$에서 시작해요. 그리고 $\|\mathbf{r}^{(t+1)}-\mathbf{r}^{(t)}\|_1 < \varepsilon$이 되거나 $I_{\max}$번 반복하면 멈춰요. $V$는 남긴 화살표들의 끝점 모음이고, 마지막 벡터는 $V$ 위에서 합이 1이 되도록 맞추어요(정규화). 지갑들은 점수가 높은 순서로도 늘어놓고, 점수가 같으면 주소로 순서를 정해요. 하지만 이런 등수는 목록을 보여 줄 때만 써요. \ref{sec:statistical-tests}절의 상관(두 값이 함께 움직이는 정도)은 원래 점수로 계산해요. 그래서 거기서는 같은 점수가 그대로 동점으로 남아요.

감쇠 계수 $d$는 순수한 퍼뜨림(전파, $d\to 1$)과 고른 점수($d\to 0$) 사이를 이어 줘요. 기본값 $0.85$는 고전적인 웹 설정과 AWP의 설정과 같아요\cend{3}. 강건성 점검(조건을 바꿔도 결과가 버티는지 보기)에서는 화살표 무게를 그대로 두고 이 값을 바꾸어 봐요(\ref{sec:scaling-robustness}절). 엄격한 $\varepsilon=10^{-8}$과 $I_{\max}=300$ 덕분에, 잰 실행 시간은 계산을 일찍 멈춘 것이 아니라 그래프의 구조에 따라 정해져요.

\dissertationSubheading[sec:data-sources]{자료 출처와 표본 뽑기}

이 절에서는 주 표본의 체인(블록체인), 기간, 공개 표, 이벤트(계약이 남기는 기록 하나하나), 지갑을 정해요. 그 뒤의 어떤 단계도 이 표본 밖의 이벤트를 끌어오지 않아요. 그림~\ref{fig:data-overview}에서는 이 논문에서 쓴 온체인(블록체인 위의) 기록을 종류마다 하나의 달력 위에 놓아 보여 줘요. 분석마다 점수를 매기는 달, 라벨로 떼어 두는 달, 그리고 각 기록을 어디에 쓰는지도 함께 보여 줘요.

\begin{figure}[htbp]
\centering
\begin{adjustbox}{max width=\linewidth}
% !TEX root = ../main.tex
% Blockchain data used in the thesis: one lane per analysis, one bar per kind of
% on-chain record, placed on the months it covers (1 Dec 2025 to 30 Sep 2026).
\begin{tikzpicture}[x=1cm, y=1cm, font=\scriptsize, oneline/.style={text height=1.6ex, text depth=0.45ex}]
\definecolor{dsApproval}{RGB}{0,114,178}
\definecolor{dsTransfer}{RGB}{0,158,115}
\definecolor{dsGMX}{RGB}{230,159,0}
\definecolor{dsAave}{RGB}{204,121,167}
\definecolor{dsOther}{RGB}{140,140,140}
\def\xo{1.95}% x of 1 Dec 2025
\def\mw{0.85}% width of one month
\def\bh{0.36}% bar height
\def\xr{10.7}% left edge of the right-hand column
% \dsInput{colour}{y}{first month}{end month}{text}: records used as input
% (scores, cohorts, same-window checks). Months count from 0 = December 2025.
\def\dsInput#1#2#3#4#5{%
  \fill[#1!35] ({\xo+\mw*(#3)},{#2-\bh/2}) rectangle ({\xo+\mw*(#4)},{#2+\bh/2});
  \draw[#1!75!black, line width=0.5pt] ({\xo+\mw*(#3)},{#2-\bh/2}) rectangle ({\xo+\mw*(#4)},{#2+\bh/2});
  \node[oneline, inner sep=0pt] at ({\xo+\mw*((#3)+(#4))/2},#2) {#5};}
% \dsHeld{colour}{y}{first month}{end month}{text}{border style}: labels counted
% after a freeze date and kept out of the scores.
\def\dsHeld#1#2#3#4#5#6{%
  \fill[pattern=north east lines, pattern color=#1] ({\xo+\mw*(#3)},{#2-\bh/2}) rectangle ({\xo+\mw*(#4)},{#2+\bh/2});
  \draw[#1!75!black, line width=0.7pt, #6] ({\xo+\mw*(#3)},{#2-\bh/2}) rectangle ({\xo+\mw*(#4)},{#2+\bh/2});
  \node[oneline, fill=white, inner sep=1pt] at ({\xo+\mw*((#3)+(#4))/2},#2) {#5};}
\def\rowlabel#1#2#3{\node[oneline, anchor=east, text=#1!75!black] at ({\xo-0.1},#2) {#3};}
\def\lanetitle#1#2{\node[oneline, anchor=west, font=\footnotesize\bfseries, fill=white, inner sep=1pt] at (0,#1) {#2};}
\def\laneuse#1#2{\node[anchor=north west, text width=5.7cm, align=flush left, inner sep=0pt] at (\xr,#1) {#2};}
\def\lanerule#1{\draw[black!25] (0,#1) -- (16.45,#1);}
\def\top{-0.3}% top of the lanes
\def\bottom{-12.8}% bottom of the lanes

% Month grid and calendar
\foreach \m in {0,...,10} {\draw[black!12] ({\xo+\mw*\m},\top) -- ({\xo+\mw*\m},\bottom);}
\foreach \m/\name in {0/12월,1/1월,2/2월,3/3월,4/4월,5/5월,6/6월,7/7월,8/8월,9/9월}
  {\node[oneline] at ({\xo+\mw*(\m+0.5)},0) {\name};}
\node[oneline, font=\scriptsize\bfseries] at ({\xo+\mw*0.5},0.34) {2025};
\node[oneline, font=\scriptsize\bfseries] at ({\xo+\mw*1.5},0.34) {2026};
\draw[decorate, decoration={brace, amplitude=4pt}] ({\xo+0.03},0.6) -- ({\xo+\mw*6-0.03},0.6)
  node[oneline, midway, above=4pt] {관찰 기간};
\draw[decorate, decoration={brace, amplitude=4pt}] ({\xo+\mw*6+0.03},0.6) -- ({\xo+\mw*9-0.03},0.6)
  node[oneline, midway, above=4pt] {등록한 라벨 기간};

% Freeze dates
\draw[black!70, dashed, line width=0.6pt] ({\xo+\mw*3},\top) -- ({\xo+\mw*3},\bottom);
\draw[black!70, dashed, line width=0.6pt] ({\xo+\mw*6},\top) -- ({\xo+\mw*6},\bottom);
\node[align=center, anchor=north] at ({\xo+\mw*3},{\bottom-0.05}) {2026년 2월 28일\\봄 동결일 $t_1$};
\node[align=center, anchor=north] at ({\xo+\mw*6},{\bottom-0.05}) {2026년 5월 31일\\등록한 동결일};

% Lane A: cohort selection
\lanetitle{-0.6}{매칭 코호트 고르기}
\rowlabel{dsGMX}{-1.0}{GMX V2}
\dsInput{dsGMX}{-1.0}{0}{6}{닫기: 지갑마다 세 번 이상}
\laneuse{-0.5}{허락이나 송금을 하나라도 꺼내기 전에 매칭 코호트(지갑 $5{,}521$개)를 정해요(\ref{sec:data-sources}절).}
\lanerule{-1.35}

% Lane B: same-window analysis
\lanetitle{-1.65}{같은 기간 분석, 매칭 코호트($n=5{,}521$)}
\rowlabel{dsApproval}{-2.05}{허락}
\dsInput{dsApproval}{-2.05}{0}{6}{엔도스랭크; 허락 점검}
\rowlabel{dsTransfer}{-2.49}{송금}
\dsInput{dsTransfer}{-2.49}{0}{6}{AWP; 송금 점검과 시빌 점검}
\rowlabel{dsGMX}{-2.93}{GMX V2}
\dsInput{dsGMX}{-2.93}{0}{6}{보관한 거래 가족}
\rowlabel{dsAave}{-3.37}{Aave V3}
\dsInput{dsAave}{-3.37}{0}{6}{보관한 대출 기준선}
\laneuse{-1.55}{허락 화살표와 송금 화살표를 함께 쓰면, 원래 금액으로 무게를 단 \mbox{C-PR}과 \mbox{S-PR}이 나와요. RQ1--RQ3, RQ5의 같은 기간 부분, 실행 시간 벤치마크(\ref{sec:benchmark-results}--\ref{sec:coupled-results}절). GMX와 Aave 행은 어떤 주장에도 쓰지 않아요(부록~\ref{ch:appendix-archive}).}
\lanerule{-3.72}

% Lane C: spring holdout on spenders
\lanetitle{-4.02}{봄 홀드아웃, 스펜더($n=1{,}335$)}
\rowlabel{dsApproval}{-4.42}{허락}
\dsInput{dsApproval}{-4.42}{0}{3}{점수 입력}
\dsHeld{dsApproval}{-4.42}{3}{6}{새 허락 쌍}{solid}
\rowlabel{dsTransfer}{-4.86}{송금}
\dsInput{dsTransfer}{-4.86}{0}{3}{점수 입력}
\dsHeld{dsTransfer}{-4.86}{3}{6}{새 송금자}{solid}
\laneuse{-3.92}{$t_1$에 허락을 받아 둔 스펜더. 점수와 차수 기준선은 $t_1$까지의 이벤트를 쓰고, 두 라벨은 그 뒤의 이벤트를 세요. RQ4, RQ5(\ref{sec:holdout-protocol}절, \ref{sec:holdout-results}절).}
\lanerule{-5.21}

% Lane D: spring trader run
\lanetitle{-5.51}{봄 트레이더 실행, 탐색적}
\rowlabel{dsApproval}{-5.91}{허락}
\dsInput{dsApproval}{-5.91}{0}{3}{점수 입력}
\rowlabel{dsTransfer}{-6.35}{송금}
\dsInput{dsTransfer}{-6.35}{0}{3}{점수 입력}
\rowlabel{dsGMX}{-6.79}{GMX V2}
\dsHeld{dsGMX}{-6.79}{3}{6}{성공 라벨}{solid}
\rowlabel{dsAave}{-7.23}{Aave V3}
\dsHeld{dsAave}{-7.23}{3}{6}{청산}{solid}
\laneuse{-5.41}{$t_1$ 그래프에 있는 매칭 지갑 4,227개예요. 그 가운데 1,860개는 3월부터 5월까지 닫기를 세 번 이상 했어요. 엔도스랭크, AWP, C-PR은 라벨을 본 뒤에, 정해 둔 규칙 없이 점수를 매겼어요(\ref{sub:trader-holdout}절, \ref{sub:trader-results}절).}
\lanerule{-7.58}

% Lane E: registered replication
\lanetitle{-7.88}{등록 재현, 2026년 6--8월}
\rowlabel{dsApproval}{-8.28}{허락}
\dsInput{dsApproval}{-8.28}{0}{6}{점수 입력}
\dsHeld{dsApproval}{-8.28}{6}{9}{새 허락 쌍}{solid}
\rowlabel{dsTransfer}{-8.72}{송금}
\dsInput{dsTransfer}{-8.72}{0}{6}{점수 입력}
\dsHeld{dsTransfer}{-8.72}{6}{9}{새 송금자}{solid}
\rowlabel{dsGMX}{-9.16}{GMX V2}
\dsHeld{dsGMX}{-9.16}{6}{9}{청산 없이 닫기}{solid}
\laneuse{-7.78}{5월 31일의 스펜더 1,964개와 5월 31일 그래프의 매칭 트레이더 5,151개(그중 \mbox{1,402개}에 라벨): F1--F6과 보고한 민감도 분석 두 가지. 등록을 커밋한 뒤 2026년 10월 1일에 꺼냈어요(\ref{sub:fresh-holdout}절, \ref{sub:fresh-results}절).}
\lanerule{-9.51}

% Lane F: runtime pool
\lanetitle{-9.81}{실행 시간 규모 변화 풀}
\rowlabel{dsOther}{-10.21}{주소}
\dsInput{dsOther}{-10.21}{0}{6}{뽑은 주소 93,559개, 화살표 없음}
\laneuse{-9.71}{확장 풀, $N=99{,}080$: 점 개수에 따른 실행 시간 민감도와 표본 정의에 따른 민감도(\ref{sub:runtime-scaling}절, \ref{sub:sample-size}절).}
\lanerule{-10.87}

% Lane G: spender account check
\lanetitle{-11.17}{스펜더 계정 점검}
\rowlabel{dsOther}{-11.57}{계정}
\dsInput{dsOther}{-11.57}{0}{3}{로그, 거래}
\fill[dsOther!75!black] ({\xo+\mw*10-0.15},-11.57) +(0,0.13) -- +(0.13,0) -- +(0,-0.13) -- +(-0.13,0) -- cycle;
\node[anchor=east, fill=white, inner sep=1pt] at ({\xo+\mw*10-0.33},-11.57) {2026년 9월 30일에 코드 읽음};
\laneuse{-11.07}{$t_1$까지 어떤 스펜더가 로그를 남기거나 거래를 보냈는지, 그다음 블록 510{,}408{,}687에서 본 스펜더마다의 코드: 1,335개 가운데 계약 \mbox{1,174개}, 외부 소유 계정 161개(\ref{sub:cohorts}절).}

% Legend: record colours, then bar styles
\begin{scope}[shift={(0,-13.85)}]
  \node[oneline, anchor=west, font=\scriptsize\bfseries, inner sep=0pt] at (0,0) {기록};
  \foreach \c/\t/\x/\y in {dsApproval/ERC-20 허락 로그/1.4/0, dsTransfer/ERC-20 송금 로그/5.0/0, dsGMX/GMX V2 포지션 이벤트/8.6/0, dsAave/Aave V3 대출 이벤트/12.45/0, dsOther/그 밖의 BigQuery 읽기와 노드 읽기/1.4/-0.42} {
    \fill[\c!35] (\x,{\y-0.1}) rectangle ({\x+0.45},{\y+0.1});
    \draw[\c!75!black, line width=0.5pt] (\x,{\y-0.1}) rectangle ({\x+0.45},{\y+0.1});
    \node[oneline, anchor=west] at ({\x+0.5},\y) {\t};
  }
  \node[oneline, anchor=west, font=\scriptsize\bfseries, inner sep=0pt] at (0,-0.9) {막대};
  \fill[dsApproval!35] (1.4,-1.0) rectangle (1.625,-0.8);
  \fill[dsTransfer!35] (1.625,-1.0) rectangle (1.85,-0.8);
  \draw[black!60, line width=0.5pt] (1.4,-1.0) rectangle (1.85,-0.8);
  \node[oneline, anchor=west] at (1.9,-0.9) {색칠: 점수, 코호트, 같은 기간 점검에 넣는 입력};
  \fill[pattern=north east lines, pattern color=black!70] (8.6,-1.0) rectangle (9.05,-0.8);
  \draw[black!60, line width=0.7pt] (8.6,-1.0) rectangle (9.05,-0.8);
  \node[oneline, anchor=west] at (9.1,-0.9) {빗금: 떼어 둔 라벨, 동결일 뒤에 셈};
  \draw[black!70, dashed, line width=0.6pt] (1.625,-1.5) -- (1.625,-1.14);
  \node[oneline, anchor=west] at (1.9,-1.32) {점선: 동결일};
\end{scope}
\end{tikzpicture}

\end{adjustbox}
\caption[{이 논문에서 쓴 블록체인 자료: 출처, 다루는 달, 떼어 둔 라벨, 쓰임새.}]{이 논문에서 쓴 블록체인 자료예요. 모두 아비트럼 원에서 왔어요. 가로줄 하나는 분석 하나이고, 막대 하나는 온체인 기록 한 종류를 그 기록이 다루는 달 위에 놓은 것이에요. 꽉 찬 막대는 점수, 코호트, 같은 기간 점검에 넣는 입력이에요. 빗금 친 막대는 동결일 뒤에 센 라벨이고, 점수에는 넣지 않았어요. 6월부터 8월까지의 로그는 재현을 등록한 뒤인 2026년 10월 1일에 꺼냈어요. 막대 안의 글은 그 기록으로 무엇을 계산하는지 알려 줘요. 오른쪽 열에는 코호트, 연구 질문, 결과를 보고하는 절이 적혀 있어요. GMX 기록으로 매칭 코호트를 고르고, 허락 질의와 송금 질의는 그 코호트에 맞추어 걸러요.}
\label{fig:data-overview}
\end{figure}

\begin{itemize}
\item 주 체인: 모든 이벤트는 아비트럼 원(체인 ID $42161$)에서 왔어요. 아비트럼 원은 이더리움 위에서 돌아가는 낙관적 롤업(거래를 일단 옳다고 보고 처리하는 롤업)이에요\cend{45}. 그 로그는 빅쿼리(BigQuery)의 공개 자료에 들어 있어요. 체인을 하나만 쓰면 체인 사이의 별칭 문제를 피할 수 있어요. 별칭 문제는 한 사용자가 여러 체인에 가진 주소들이 서로 다른 점처럼 나타나는 일이에요. 또 표본의 모든 지갑이 똑같은 수수료와 똑같은 지연 시간(처리되기까지 기다리는 시간)을 겪어요.

\item 관찰 기간: 2025-12-01 00:00:00 UTC부터 2026-05-31 23:59:59 UTC까지 여섯 달이에요. 질의에서는 2026-06-01을 끝으로 잡되, 그 시각 자체는 넣지 않아요. 여섯 달이면 활동 기간(tenure) 대리 지표와 활동한 달 수 대리 지표가 지갑마다 충분히 달라질 수 있어요. 두 그래프는 정확히 같은 날짜들을 써요. 자료 꺼내기는 질의 한 번에 쓸 수 있는 바이트 예산 안에 머물도록 한 달씩 해요(\ref{sec:collection-strategy}절). 시간 감쇠의 기준 시점은 $T_{\mathrm{obs}}=$ 2026-05-31 23:59:59 UTC예요. 그래서 마지막 날의 이벤트는 $0\le\Delta t<1$이고, 첫날의 이벤트는 약 $182$일 묵은 것이에요.

\item 빅쿼리 출처: 로그는 빅쿼리의 다음 공개 표에서 읽어요.
\begin{center}
\texttt{bigquery-public-data.goog\_blockchain\_arbitrum\_one\_us.logs}.
\end{center}
기록 하나는 계약(스마트 계약, 블록체인 위에서 저절로 실행되는 프로그램)이 내보낸 이벤트 하나예요. 기록에는 블록 시각, 주소, 토픽, 데이터가 들어 있어요. 질의는 블록 시각으로 거르고, 필요하면 계약 주소로도 거르고, 이벤트 토픽 해시로도 걸러요. 구글 클라우드(Google Cloud) 계정이 있는 사람은 누구나 같은 표에 같은 질의를 다시 돌릴 수 있어요. 부록~\ref{ch:appendix-sql}에 질의 모양을 실었어요.

\item 평판 이벤트: 엔도스랭크는 ERC-20 \texttt{Approval} 로그를 써요\cend{18}. 이 로그에는 EIP-2612 \texttt{permit}(서명 허락)으로 한 허락도 들어 있어요\cend{19}. AWP는 \texttt{Transfer} 로그를 써요. 표본 지갑이 주인이나 스펜더, 또는 보낸 쪽이나 받는 쪽인 이벤트만 남겨요. 이렇게 하면 꺼내는 자료의 양이 한도 안에 머물고, 코호트 둘레의 부분 그래프만 따로 떼어 낼 수 있어요.

\item 매칭 코호트: 일치도는 아비트럼 원의 GMX~V2 지갑 $n=5{,}521$개에 대해 재요. 이 지갑들은 2025-12-01과 2026-05-31 사이에 (포지션) 닫기를 세 번 이상 했어요. 여기서 닫기 한 번은 GMX \texttt{PositionDecrease} 이벤트 하나예요. 그래서 일부만 닫은 것과 청산(손실이 너무 커져 거래가 강제로 끝나는 일)도 닫기로 세어요. 이 목록은 GMX 이벤트 자료만으로 만들고, 허락이나 송금을 하나라도 꺼내기 전에 고정해요. 그래서 코호트는 평판 그래프에 따라 달라지지 않고, 점수 사이의 차이가 표본 뽑기 때문에 생길 수 없어요. 코호트 지갑이 꼭 어느 그래프에 나와야 하는 것은 아니에요. $966$개에는 허락 화살표가 하나도 닿지 않고, $381$개에는 송금 화살표가 하나도 닿지 않고, $370$개에는 둘 다 닿지 않아요. 그런 지갑은 해당 방법에서 점수가 $0$이에요(\ref{sec:graph-construction}절).

\item 빼는 것: 멤풀(mempool, 아직 블록에 들어가지 않은 거래가 기다리는 곳)의 거래, 오프체인(블록체인 밖의) 주문장, 중앙화 거래소 계정, 브리지(체인과 체인을 잇는 다리)는 자료에 들어 있지 않아요. 또 한 운영자가 움직이는 여러 주소를 한 주체로 합치려고 하지도 않아요.

\end{itemize}

\dissertationSubheading[sec:collection-strategy]{두 단계 수집 전략}

자료 모으기와 평가는 서로 떨어진 두 단계로 해요. 그래서 우리든, 자료 꺼내기를 해 본 다른 사람이든, 표와 벤치마크와 강건성 실행을 클라우드 비용 없이 필요한 만큼 몇 번이고 다시 만들 수 있어요.

\begin{itemize}
\item 1단계(평판 자료 흐름): GMX 자료로 코호트를 고정한 뒤, 코호트 지갑이 주인이나 스펜더인 \texttt{Approval} 이벤트를 모두 꺼내고, 코호트 지갑이 보낸 쪽이나 받는 쪽인 \texttt{Transfer} 이벤트도 모두 꺼내요. 질의는 한 달에 하나씩 하고, 질의마다 미리 돌려 보기(dry run)로 비용을 따로 어림해요. 허락은 (블록 번호, 로그 번호) 순서로 보아, 주인--스펜더--토큰 세 값 묶음마다 가장 최근의 허락 하나로 줄여요. 허락이 $0$으로 보이면 그 화살표를 버려요. 송금 양은 원래 단위 그대로 두고, 양이 0인 송금은 버려요. 그래프를 만드는 데 들어가는 입력은 이 표들뿐이에요.

\item 2단계(로컬 평가): 1단계에서 만든 parquet 파일로, 로컬 컴퓨터(연구자 자신의 컴퓨터)가 점수, 점검, 일치도 통계, 벤치마크, 강건성 실행을 계산해요. 확장 지갑 풀은 평판 부분 그래프에 있는 주소와 보충 parquet 파일에 있는 주소로 만들어요. 그래서 체인을 다시 한 번 훑을 필요가 없어요.

\item 비용 안전장치: 질의 한 번의 예산은 $150$\,GiB예요\footnote{질의마다 바이트 양을 이렇게 제한하는 것은 이 연구의 처리 과정이 스스로 하는 일이에요. 빅쿼리가 정한 제한은 아니에요.}. 미리 돌려 보기로 어림한 양이 이 예산을 넘는 질의는, 운영자가 추가 플래그(따로 붙이는 확인 표시)로 확인하지 않으면 멈춰요. 질의마다 훑은 지갑, 돌려받은 행, 처리한 바이트, 질의 시간을 처리 과정과 함께 기록해요(부록~\ref{ch:appendix-eval}). 그래서 자료를 얻는 데 든 비용을 따져 볼 수 있어요.

\end{itemize}

\dissertationSubheading[sec:pipeline-architecture]{처리 과정의 구조}

그림~\ref{fig:pipeline-architecture}에서는 공개 로그에서 연구 표까지, 처리 과정 전체를 따라가며 보여 줘요. 원격 단계(클라우드에서 하는 단계)는 모두 바꿀 수 없는 parquet 파일을 써요. 입력과 버전을 매긴 설정 파일이 정해지면, 모든 로컬 단계는 언제나 같은 결과를 내요.

\begin{figure}[htbp]
\centering
\begin{adjustbox}{max width=\linewidth}
\begin{tikzpicture}[
  font=\small,
  node distance=1.0cm and 1.0cm,
  box/.style={draw, rounded corners=2pt, align=center, minimum height=0.95cm, minimum width=2.2cm, inner sep=3pt},
  wide/.style={draw, rounded corners=2pt, align=center, minimum height=0.95cm, minimum width=2.6cm, inner sep=3pt},
  arrow/.style={-{Stealth[length=2.5mm]}, thick}
]
% Top row: remote acquisition (left to right)
\node[wide] (bq) at (0,0) {BigQuery\\공개 로그};
\node[wide] (extract) at (3.2,0) {로그 꺼내기\\(달마다; 토픽과\\지갑으로 거르기)};

% Local evaluation. Row 1: events -> edge lists -> scores -> evaluation -> export.
% Row 2: proxy metrics from the same events. Row 3: temporal holdout.
\node[box] (decode) at (0,-2.3) {이벤트\\풀어 읽기};
\node[box] (prep) at (3.0,-2.3) {화살표 목록\\만들기};
\node[wide] (ranks) at (6.2,-2.3) {페이지랭크 점수\\(EndorseRank, AWP,\\C-PR, S-PR)};
\node[wide] (align) at (9.5,-2.3) {일치도\\\& 벤치마크};
\node[wide] (tabexp) at (12.7,-2.3) {표 \&\\그림 내보내기};
\node[box] (proxy) at (3.0,-3.9) {대리\\지표};
\node[wide] (hold) at (6.2,-5.2) {시간 홀드아웃\\($t_1$의 점수, $t_2$까지의 라벨)};

\draw[arrow] (bq) -- (extract);
\draw[arrow] (extract.south) -- ++(0,-0.55) -| (decode.north);
\draw[arrow] (decode) -- (prep);
\draw[arrow] (prep) -- (ranks);
\draw[arrow] (ranks) -- (align);
\draw[arrow] (align) -- (tabexp);
\draw[arrow] ([xshift=4pt]decode.south) |- (proxy.west);
\draw[arrow] ([xshift=-4pt]decode.south) |- (hold.west);
\draw[arrow] (proxy.east) -| (align.south);
\draw[arrow] (hold.east) -| (tabexp.south);

\node[draw, dashed, rounded corners, fit=(bq)(extract), inner sep=8pt, label=above:{\footnotesize 원격 / 자료 얻기}] {};
\node[draw, dashed, rounded corners, fit=(decode)(prep)(ranks)(align)(tabexp)(proxy)(hold), inner sep=8pt, label=below:{\footnotesize 로컬 / 다시 할 수 있음}] {};
\end{tikzpicture}
\end{adjustbox}
\caption[{처리 과정 한눈에 보기: 아비트럼 공개 로그를 나누어 읽고 풀어서 화살표 기록을 준비하고, 엔도스랭크와 AWP 점수를 계산하고, 일치도와 벤치마크 계산을 한 다음, 결과 표를 쓰는 과정.}]{처리 과정 한눈에 보기예요. 허락 로그와 송금 로그는 원격에서 한 달 단위 덩어리로 걸러서 꺼내요. 로컬에서는 이벤트를 풀어 읽어 화살표 목록으로 바꾸고, 그 목록으로 엔도스랭크, AWP, C-PR, S-PR을 계산해요. 대리 지표는 나란히 놓인 다른 갈래에서 같은 이벤트로 만들어요. 그다음 일치도와 벤치마크 계산, 시간 홀드아웃이 이어지고, 마지막으로 결과 표와 그림을 내보내요.}
\label{fig:pipeline-architecture}
\end{figure}

1단계 스크립트는 달마다 나눈 조각 파일과, 이를 모아 정리한 마지막 허락 표와 송금 표를 써요. 2단계 스크립트는 이것들을 읽어 그래프를 만들고, 점수를 계산하고, \ref{ch:results}장의 표를 만들어요. 처리 과정 전체는 시험용으로 만든 가짜 자료(합성 고정 자료)로 인터넷 없이도 돌릴 수 있어요.

\dissertationSubheading[sec:event-decoding]{이벤트 풀어 읽기}

스마트 계약이 내보낸 로그에는 토픽과 데이터 내용이 있지만, 칸 이름은 없어요. 풀어 읽기(해독)는 각 행을 그래프 만들기와 점검에 쓸, 종류가 정해진 기록으로 바꾸는 일이에요. 여기서 실수가 나면 두 그래프가 아무 표시 없이 조용히 망가져요. 토픽 해시와 칸 대응은 아래에 적었어요.

\begin{itemize}
\item 로그의 생김새: 로그마다 그것을 내보낸 계약, 토픽 여러 개, ABI(계약과 자료를 주고받는 규칙) 방식으로 묶은 데이터 덩어리가 기록돼요. 첫 번째 토픽은 이벤트 서명(이벤트 이름과 칸 종류를 적은 글)의 keccak-256 해시예요. 색인된 매개변수(따로 찾을 수 있게 표시한 값)는 그 뒤에 $32$바이트 낱말로 이어지고(주소는 왼쪽을 0으로 채워요), 색인되지 않은 매개변수는 ABI 순서대로 데이터 칸에 들어가요\cend{18}. ERC-721 토큰도 같은 서명으로 \texttt{Transfer}와 \texttt{Approval}을 내보내요. 하지만 토큰 ID를 네 번째 토픽에 넣고 데이터는 비워 둬요. 그래서 그 로그들은 0으로 풀리고, 양이 0인 다른 기록들과 함께 버려져요. SQL은 토픽으로 거르고, Python은 숫자 칸을 풀어 읽어요.

\item ERC-20 \texttt{Approval}: 이 이벤트의 토픽 해시는
\begin{center}
\texttt{0x8c5be1e5ebec7d5bd14f71427d1e84f3dd0314c0f7b2291e5b200ac8c7c3b925},
\end{center}
곧 \texttt{Approval(address,address,uint256)}의 keccak-256 해시예요. 주인과 스펜더는 토픽~$1$과 토픽~$2$에 들어 있어요. 허락한 양은 데이터 칸에서 토큰의 기본 단위로 풀어 읽고, 소수 자릿수(decimals)로 나누어 맞추지 않아요. 무제한 허락(\texttt{type(uint256).max})은 원래 값을 그대로 가져요. \texttt{permit}도 같은 \texttt{Approval} 이벤트를 내보내요. 그래서 허락은 온체인에서 \texttt{approve()}를 불러서 했든, 오프체인 서명으로 했든 똑같이 세어요\cend{19}.

\item ERC-20 \texttt{Transfer}: 이 이벤트의 토픽 해시는
\begin{center}
\texttt{0xddf252ad1be2c89b69c2b068fc378daa952ba7f163c4a11628f55a4df523b3ef},
\end{center}
곧 \texttt{Transfer(address,address,uint256)}의 keccak-256 해시예요. 보낸 쪽과 받는 쪽은 색인되어 있고, 양은 기본 단위 그대로 둬요. 0 주소를 쓰는 발행(mint, 새 토큰 만들기) 로그와 소각(burn, 토큰 없애기) 로그는 코호트 지갑이 관련되어 있으면 언제나 남겨 둬요. 자기에게 보낸 송금은 이벤트 표에 남아 있지만, 들어오는 상대를 세는 계산에서는 빼요(\ref{sec:proxy-formulas}절). 그런 송금으로는 상대가 누구인지 알 수 없기 때문이에요.

\item 주소 맞추기와 단위: 주소는 표를 이어 붙이기(join) 전에 모두 소문자로 바꿔요. 양은 \texttt{uint256} 값을 원래 단위 그대로 배정밀도 실수(소수점이 있는 수를 담는 컴퓨터 형식)로 담은 것이에요. 소수 자릿수로 나누지도, 달러(USD)로 바꾸지도, 위쪽을 자르거나 극단값을 줄이지도(윈저화) 않아요. 그래서 무제한 허락은 $2^{256}-1\approx1.16\times10^{77}$로 들어가요. 엔도스랭크에서는 상한이 있는 변환 $V$ 때문에 이 허락도 다른 허락처럼 세어요(\ref{sub:endorserank-graph}절). C-PR의 금액 그대로인 허락 층에서는 이 허락이 그 주인의 행을 거의 다 차지해요(\ref{sub:coupled-graph}절).

\end{itemize}

\dissertationSubheading[sec:graph-construction]{그래프 만들기}

엔도스랭크와 AWP는 화살표에 무게를 다는 방법이 같고, 같은 페이지랭크 풀이 프로그램(식~\eqref{eq:pagerank})을 돌려요. 다른 점은 화살표를 만드는 이벤트와, 걷기가 어디서 다시 출발하는지예요. 두 방법 모두 끝점 가운데 적어도 하나가 씨앗 지갑 모음 $\mathcal{W}$에 있는 화살표만 남겨요. 그러면 씨앗 지갑에서 한 걸음 떨어진 이웃들도 점으로 들어와요. 그리고 각 방법은 $\mathcal{W}$의 모든 지갑에 점수를 보고해요. 부분 그래프에 없는 지갑은 점수가 $0$이에요.

엔도스랭크의 화살표는 보증(믿고 맡김)을 나타내고, AWP의 화살표는 송금을 나타내요. 허락은 주인 대신 토큰을 써도 된다고 일부러 내주는 권한이고, 나중에 거둘 수도 있어요. 송금은 이미 끝난 값의 이동이에요. 송금은 돈 내기, 여러 곳을 거쳐 가는 길(라우팅), 시장 조성(늘 사고팔 값을 내놓아 거래를 돕는 일), 자전 거래(자기끼리 사고팔아 거래가 많은 척하기)의 한 부분일 수 있어요. 이 가운데 어느 것도 보낸 쪽이 보증을 약속하는 것은 아니에요\cend{4,18}. 두 방법 모두 이벤트에 나이에 따른 무게를 달아서, 최근 이벤트를 더 무겁게 세요. 또 양을 상한이 있는 함수에 넣어서, 이벤트 하나가 많아야 한 번으로만 세지게 해요. 엔도스랭크가 AWP의 무게를 가져다 쓰므로, AWP를 먼저 설명해요.

\dissertationSubsubheading[sub:awp-graph]{AWP 송금 그래프}

AWP는 Do, Do, Nguyen이 IEEE RIVF 2023에서 발표한 방법이에요\cend{3}. 토큰 양이 $x_e$이고, 관찰 기간이 끝나는 $T_{\mathrm{obs}}$보다 $\Delta t_e$일 전에 일어난 송금 $e$를 생각해 봐요. 이 송금은 보낸 쪽에서 받는 쪽으로 가는 화살표에 $\sigma(\Delta t_e)\,V(x_e)$의 무게로 들어가요. 시간 무게는 다음과 같은 로지스틱 감쇠예요.
\begin{equation}
\label{eq:logistic-decay}
\sigma(\Delta t)
=
\frac{1}{1+\exp\bigl(k(\Delta t-t_0)\bigr)},
\qquad
k=0.01,\quad t_0=180\text{일},
\end{equation}
값 무게는 다음과 같아요.
\begin{equation}
\label{eq:value-transform}
V(z)
=
\frac{2}{1+e^{-bz}}-1,
\qquad
b=1,
\end{equation}
이 값은 $z=0$에서 $0$이고, 양이 커질수록 $1$을 향해 올라가요. 화살표 무게는 송금들을 더한 값이에요.
\begin{equation}
\label{eq:transfer-weight}
w_T(u,v)
=
\sum_{e:\,u\to v}
\sigma(\Delta t_e)\,V(x_e).
\end{equation}
이 합은 토큰이 무엇이든, 기간 안에 $u$에서 $v$로 간 모든 송금을 더해요. 자기에게 보낸 송금($u=v$)도 남기고, 이것은 $u$에서 자기에게 돌아오는 고리(자기 고리)가 돼요. 걷기는 주소의 활발함에 비례해서 그 주소에서 다시 출발해요.
\begin{equation}
\label{eq:activeness}
X_u
=
\sum_{v\neq u}\;
\max_{e:\,u\to v}\sigma(\Delta t_e),
\end{equation}
이 식은 $u$가 송금을 보낸 주소마다, 그 주소로 보낸 마지막 송금의 시간 무게를 더한 거예요. 식~\eqref{eq:pagerank}의 재시작 분포는 $\pi_u=X_u/\sum_{v\in V_T}X_v$이고, 이 분포가 매달린 점들의 몫도 받아요. 다른 주소에 한 번도 보내지 않은 주소는 $X_u=0$이라서 재시작 몫을 받지 못해요. 그래서 그 주소의 점수는 들어오는 화살표에서만 나와요. 감쇠 계수는 $d=0.85$예요.

그 논문은 이더리움 거래에 점수를 매겼고, 매개변수 값은 알려 주지 않아요. 그 논문의 시간 무게는 식~\eqref{eq:logistic-decay}의 로지스틱 함수를 나이 대신 달력 시간으로 쓴 것이에요. $k$, $t_0$, $b$는 이 연구가 고른 값이에요. 양을 원래 단위로 두면(\ref{sec:event-decoding}절), 기본 단위 $10$개 이상인 모든 송금에서 $V(x_e)$와 $1$의 차이가 $10^{-4}$ 안이에요. 이런 송금이 기간 안 송금의 $99.5\%$예요. 그래서 화살표 무게는 사실상 그 화살표 위의 송금 수이고, 송금 하나하나를 그 시간 무게만큼 세요. 주소의 활발함도 사실상 그 주소가 송금을 보낸 주소의 수를 같은 방식으로 센 것이에요. 기간 안의 나이에서는 감쇠가 거의 직선이에요(그림~\ref{fig:logistic-decay-methods}). 무게는 $\sigma(0)=0.858$에서 $\sigma(90)=0.711$을 거쳐 $\sigma(182)=0.495$로 떨어져요. 그리고 $t_0$보다 오래된 이벤트는 기간의 처음 이틀 동안 일어난 이벤트뿐이에요.

\begin{figure}[htbp]
\centering
\includegraphics[width=0.8\textwidth]{\figpath{logistic-decay.pdf}}
\caption[{AWP와 엔도스랭크의 로지스틱 시간 감쇠 $\sigma(\Delta t)$ ($k=0.01$, $t_0=180$일). 가로축은 이벤트가 2026-05-31 UTC보다 며칠 전에 일어났는지(이벤트의 나이), 세로축은 송금이나 허락의 시간 무게.}]{AWP와 엔도스랭크의 로지스틱 시간 감쇠 $\sigma(\Delta t)$예요($k=0.01$, $t_0=180$일). 가로축은 이벤트가 2026-05-31 UTC보다 며칠 전에 일어났는지(이벤트의 나이)를 보여 줘요. 세로축은 송금이나 허락의 시간 무게를 보여 줘요. 색칠한 띠는 관찰 기간 안에 나오는 나이(0일부터 182일까지)를 표시해요. 이 범위에서 무게는 $0.86$에서 $0.50$까지만 떨어져요.}
\label{fig:logistic-decay-methods}
\end{figure}

알고리즘~\ref{alg:awp}에 그래프를 만들고 점수를 매기는 단계를 정식으로 적었어요.

\begin{algorithm}[htbp]
\caption{AWP 그래프 만들기와 점수 매기기}
\label{alg:awp}
\begin{algorithmic}[1]
\Require 송금들 $\mathcal{T}$, 씨앗 지갑들 $\mathcal{W}$, 관찰 끝 시각 $T_{\mathrm{obs}}$, 매개변수 $k,t_0,b,d,\varepsilon,I_{\max}$
\Ensure $\mathcal{W}$ 위의 점수 $\mathbf{s}_{\mathrm{AWP}}$
\State $E \gets \emptyset$; 모든 주소 $u$에 대해 $X_u\gets 0$
\For{$\mathcal{T}$의 각 송금 $e=(u,v,x_e,t_e)$}
  \State $\Delta t \gets (T_{\mathrm{obs}}-t_e)$ (일 단위)
  \State 화살표 $(u\to v)$에 $\sigma(\Delta t)\,V(x_e)$를 더해 쌓아요 \Comment{식~\eqref{eq:logistic-decay}--\eqref{eq:transfer-weight}; $u=v$이면 자기 고리가 돼요}
\EndFor
\For{송금을 보낸 각 주소 $u$}
  \State $X_u \gets \sum_{v\neq u}\max_{e:\,u\to v}\sigma(\Delta t_e)$ \Comment{식~\eqref{eq:activeness}}
\EndFor
\State 무게가 0 이하인 화살표를 버리고, $\mathcal{W}$에 닿는 화살표만 남겨요
\State 남긴 화살표의 점 $u$마다 $\pi_u \gets X_u/\sum_{v}X_v$
\State $\mathbf{s} \gets \mathrm{WeightedPageRank}(E, \boldsymbol{\pi}, d, \varepsilon, I_{\max})$ \Comment{식~\eqref{eq:pagerank}--\eqref{eq:transition-impl}}
\State \Return $i\in\mathcal{W}$마다 $\mathbf{s}_{\mathrm{AWP}}[i] \gets \mathbf{s}[i]$ (없으면 $0$)
\end{algorithmic}
\end{algorithm}

\dissertationSubsubheading[sub:endorserank-graph]{엔도스랭크 허락 그래프}

엔도스랭크는 AWP의 화살표 무게를 허락 그래프에 얹어요. 미리 다듬은(전처리한) 뒤의 마지막 허락들의 모음을 $\mathcal{A}$라고 해요. (토큰, 주인--스펜더) 묶음마다, (블록 번호, 로그 번호) 순서로 보아 관찰 기간 안의 가장 최근 \texttt{Approval} 이벤트 하나만 남겨요. 허락이 0이면 그 화살표를 지워요. 살아남은 허락 $a^{\star}(u,v,\mathrm{token})$은 $T_{\mathrm{obs}}$보다 $\Delta t^{\star}(u,v,\mathrm{token})$일 전에 정해졌어요. 이 허락은 그 나이와 그 양을 가진 이벤트 하나로 화살표에 들어가요.
\begin{equation}
\label{eq:approve-weight}
w_A(u,v)
=
\sum_{\mathrm{token}}
\sigma\bigl(\Delta t^{\star}(u,v,\mathrm{token})\bigr)\,
V\bigl(a^{\star}(u,v,\mathrm{token})\bigr),
\end{equation}
여기서 $\sigma$, $V$, $k$, $t_0$, $b$는 AWP와 같아요. 방향이 있는 화살표 $u\to v$는 주인 $u$가 스펜더 $v$를 보증한다는 뜻이에요. 거의 모든 허락(허락의 $99.9\%$)은 기본 단위 $10$개보다 훨씬 커요. 그래서 $w_A(u,v)$는 사실상 $u$가 $v$에게 허락을 해 준 토큰의 수이고, 토큰 하나하나를 그 허락의 나이에 따라 세요. 무제한 허락도 작은 허락보다 더 많이 세지 않아요. 감쇠 계수는 $d=0.85$예요.

걷기는 고르게 다시 출발해요. AWP의 재시작 규칙을 허락에 그대로 옮기면, 최근에 허락을 해 준 스펜더가 많은 주인일수록 걷기가 그 주인에게서 더 자주 다시 출발하게 돼요. 그러면 보증을 해 주는 쪽이 상을 받게 돼요. 하지만 엔도스랭크는 보증을 받는 지갑에 순위를 매기려는 방법이에요. 우리는 이 연구의 자료로 두 재시작 규칙 모두 점수를 매겨 본 뒤에 이렇게 골랐어요. 그래서 AWP의 재시작을 쓴 변형도 엔도스랭크 옆에 함께 보고해요(표~\ref{tab:endorserank-restarts}).

무게는 마지막 허락의 양과 시각을 담아요. 기간이 끝날 때 실제로 남아 있는 허락 양은 아니에요. 한 쌍 안에서 토큰마다, 가장 새 허락이 예전 허락들을 대신해요\cend{18}. 2025년 12월 1일 전에 해 준 허락은 보이지 않아요. 정해진 양만큼의(유한한) 허락은 그 뒤에 \texttt{transferFrom}이 줄였더라도 허락한 값을 그대로 가져요. 그리고 Uniswap의 Permit2 같은 허락 관리 계약에 허락을 해 준 주인은, 그 계약으로 가는 화살표로만 나타나요. 나중에 그 계약을 거쳐 토큰을 쓰는 응용 프로그램으로 가는 화살표는 생기지 않아요.

그래프를 만들고 페이지랭크를 돌리는 단계는 알고리즘~\ref{alg:endorserank}에 적었어요.

\begin{algorithm}[htbp]
\caption{엔도스랭크 그래프 만들기와 점수 매기기}
\label{alg:endorserank}
\begin{algorithmic}[1]
\Require 마지막 허락들 $\mathcal{A}$, 씨앗 지갑들 $\mathcal{W}$, 관찰 끝 시각 $T_{\mathrm{obs}}$, 매개변수 $k,t_0,b,d,\varepsilon,I_{\max}$
\Ensure $\mathcal{W}$ 위의 점수 $\mathbf{s}_{\mathrm{ER}}$
\State $E \gets \emptyset$
\For{$\mathcal{A}$의 각 마지막 허락 $(u,v,\mathrm{token},a^{\star},t^{\star})$}
  \State $\Delta t^{\star} \gets (T_{\mathrm{obs}}-t^{\star})$ (일 단위)
  \State 화살표 $(u\to v)$에 $\sigma(\Delta t^{\star})\,V(a^{\star})$를 더해 쌓아요 \Comment{식~\eqref{eq:approve-weight}}
\EndFor
\State 무게가 0 이하인 화살표를 버리고, $\mathcal{W}$에 닿는 화살표만 남겨요
\State $\mathbf{s} \gets \mathrm{WeightedPageRank}(E, \text{고른 }\boldsymbol{\pi}, d, \varepsilon, I_{\max})$ \Comment{식~\eqref{eq:pagerank}--\eqref{eq:transition-impl}}
\State \Return $i\in\mathcal{W}$마다 $\mathbf{s}_{\mathrm{ER}}[i] \gets \mathbf{s}[i]$ (없으면 $0$)
\end{algorithmic}
\end{algorithm}

$G_{\text{approve}}$는 오래 이어지는 권한을 기록하고, $G_{\text{transfer}}$는 지나간 값의 흐름을 기록해요. 허락은 드물고 뜻이 분명한 이벤트이고, 송금은 흔하고 기계적으로 일어날 때가 많아요. 그래서 매칭 코호트에서는 허락 부분 그래프가 훨씬 듬성듬성해요(\ref{ch:results}장). 페이지랭크 반복 한 번에 드는 비용이 $O(|E|)$이므로, 이 희소성(화살표가 적음)이 가설 H3의 구조적 바탕이에요.

두 방법 모두 화살표 목록을 (출발, 도착, 무게) 세 값 묶음으로 저장해요. 겹치는 화살표는 무게를 더해 하나로 합치고, 무게가 0 이하인 것은 풀기 전에 버려요. 풀이 프로그램은 듬성듬성한 행 압축 형식(0이 아닌 칸만 행마다 모아 저장하는 형식)으로 이동 행렬을 만들고, 재시작 벡터도 만들어요. 그리고 식~\eqref{eq:pagerank}의 계산을 되풀이해서 점마다 점수를 돌려줘요. 코호트 점수는 $\mathcal{W}$에 맞추어 늘어놓고, 빠진 칸은 0으로 채워요.

\dissertationSubsubheading[sub:coupled-graph]{두 층 위의 결합 연산자(C-PR)와 씨앗을 쓰는 빼 보기 실험(S-PR)}

엔도스랭크와 AWP는 저마다 한 종류의 화살표 위만 걸어요. RQ5는 두 종류의 화살표를 한 번에 걷는 걷기가 한 종류만 걷는 걷기보다 지갑에 순위를 더 잘 매기는지 물어요. 결합 연산자는 2026년 9월 14일에, 이벤트마다 원래 금액 그대로를 무게로 다는 두 층 위에서 정했어요.
\begin{equation}
\label{eq:raw-layers}
\tilde w_A(u,v)
=
\sum_{\mathrm{token}}
a^{\star}(u,v,\mathrm{token}),
\qquad
\tilde w_T(u,v)
=
\sum_{e:\,u\to v}
x_e\,\sigma(\Delta t_e).
\end{equation}
허락 층에는 엔도스랭크의 화살표가 있고, 송금 층에는 AWP의 화살표가 있어요. 다른 것은 무게뿐이에요. 양은 원래 단위로 들어가고, 상한이 있는 변환도, 소수 자릿수 맞추기도, 위쪽 자르기도 하지 않아요. 그리고 허락 층에는 시간 감쇠가 없어요. RQ5에 대한 두 판정 규칙은 모두 이 층들 위에서 정했어요. 그래서 연산자를 정한 그대로 써요. 두 층은 합친 점 모음 $V=V_A\cup V_T$ 위에서 그래프 하나를 이뤄요. 이동 행렬 $P_A$와 $P_T$는 식~\eqref{eq:transition-impl}에 따라 구해요. 지갑 $u$의 결합 행은 두 층의 행을 볼록 혼합(더해서 1이 되는 비율로 섞기)한 것이에요. 이때 비율은 $u$에게 나가는 화살표가 있는 층들만 놓고 다시 맞춰요.
\begin{align}
P^{(\lambda)}_{u,\cdot}
&=
\alpha_A(u)\,P_{A,u,\cdot}+\alpha_T(u)\,P_{T,u,\cdot},
\label{eq:coupled-transition}\\
\alpha_A(u)
&=
\frac{\lambda\,\mathbf{1}[o_A(u)>0]}{D_{\lambda}(u)},
\qquad
\alpha_T(u)=\frac{(1-\lambda)\,\mathbf{1}[o_T(u)>0]}{D_{\lambda}(u)},
\nonumber
\end{align}
여기서 $D_{\lambda}(u)=\lambda\,\mathbf{1}[o_A(u)>0]+(1-\lambda)\,\mathbf{1}[o_T(u)>0]$이고, $o_A(u)$와 $o_T(u)$는 두 층에서 $u$의 나가는 무게의 합이에요. 한 층에만 나가는 화살표가 있는 지갑은 무게 1로 그 층을 따라가요. 무게가 $0$인 층은 그 화살표와 점까지 함께 빼요. $D_{\lambda}(u)=0$인 지갑은 매달린 점이고, 그 몫은 식~\eqref{eq:pagerank}의 세 번째 항이 맡아요. 순간 이동은 고르게 하고, $d=0.85$예요. $\mathbf{s}_{\mathrm{CPR}}$는 $P=P^{(\lambda)}$로 둔 식~\eqref{eq:pagerank}의 페이지랭크 벡터이고, 풀이 설정은 한 층 방법들과 같아요.

$\lambda=1$이면 연산자는 허락 층만 걷고, $\lambda=0$이면 송금 층만 걸어요. 이 두 걷기를 C-PR ($\lambda=1$)과 C-PR ($\lambda=0$)이라고 써요. 이 둘은 RQ5에 대한 두 규칙에서 모두 비교 대상이에요(\ref{sec:holdout-protocol}절과 \ref{sub:fresh-holdout}절). 등록 문서에서는 이 둘을 엔도스랭크와 AWP라고 불러요. 단위 시험(프로그램의 작은 부분을 따로 확인하는 시험)은 시험용 고정 자료에서, 양 끝 값마다 그 층만 고르게 다시 출발하며 걸은 점수가 나오는지 확인해요. 이 C-PR 묶음은 양 끝에서 이어지지 않아요(불연속이에요). $0<\lambda<1$이면 점 모음은 $V_A\cup V_T$이고, 한 층에만 나가는 화살표가 있는 점은 $\lambda$가 얼마이든 그 층을 따라가요. 그런데 $\lambda=1$에서는 나가는 화살표가 모두 송금인 점이 매달린 점이 되거나 아예 빠져요. 그래서 $\lambda$가 $1$에 가까워질 때 C-PR은 그런 점들의 송금 화살표를 여전히 따라가는 걷기에 가까워져요. 반대쪽 끝에서도 마찬가지예요. $\lambda=0.5$(주 분석) 값과 $\lambda\in\{0.25,0.75\}$(민감도 분석) 값은 2026년 9월 14일에 설정에 정해 두었어요. 결합 점수를 하나도 계산하기 전이었어요.

금액을 그대로 쓰면 허락 층이 아주 고르지 않게 돼요. 마지막 허락 $21{,}585$개 가운데 $13{,}055$개($60.5\%$)가 무제한 허락이에요. 허락 부분 그래프에서는 화살표의 $51\%$가 무제한 허락을 담고 있어요. 주인 $4{,}590$명 가운데 $3{,}986$명이 무제한 허락을 적어도 하나 가지고 있어요. 배정밀도로 계산하면, 무제한 허락 옆에 있는 유한한 허락은 그 주인의 행에서 거의 없는 것과 같은 몫만 받아요. 그래서 그런 주인의 행은 그 주인의 무제한 허락들끼리 나누어 가져요. 두 종류를 모두 가진 주인 $1{,}987$명의 유한한 허락은 거의 아무 영향도 주지 못해요. 송금 층에서는 서로 다른 토큰의 양을 각자의 원래 크기 그대로 더해요. 그래서 소수 자릿수가 많은 토큰이 적은 토큰보다 훨씬 무겁게 세져요. 엔도스랭크와 AWP는 $V$ 덕분에 이 두 효과를 모두 피해요.

S-PR 빼 보기 실험에서는 송금 층을 걸어요. 하지만 식~\eqref{eq:pagerank}의 고른 순간 이동 벡터 $\boldsymbol{\pi}$ 대신, C-PR ($\lambda=1$)의 점수를 $V_T$로 좁히고 그 위에서 합이 1이 되게 다시 맞춘 것을 써요. $\mathbf{r}_{A}$를 $V_A$ 위에서 허락 층만 걸어 얻은 페이지랭크 벡터라고 하고, $u\notin V_A$이면 $r_{A,u}=0$으로 넓혀 둬요. 그러면 다음과 같아요.
\[
\pi_u=\frac{r_{A,u}}{\sum_{v\in V_T} r_{A,v}}\qquad(u\in V_T),
\]
그래서 허락 점수가 없는 송금 그래프의 점은 재시작 몫을 하나도 받지 못해요. $V_T$의 점 가운데 $V_A$에 있는 점이 하나도 없으면 분모가 0이 되고, 그때 $\boldsymbol{\pi}$는 $V_T$ 위의 고른 분포로 돌아가요. 같은 벡터로 매달린 점의 몫도 다시 나누어요. 그 결과는 트러스트랭크(TrustRank)\cend{17}와 같은 짜임이에요. 다만 사람이 손으로 표시한 씨앗 대신 허락에서 나온 씨앗을 써요. 허락은 걷기가 어디서 다시 출발할지 정하고, 송금은 걷기가 어떻게 퍼질지 정해요. S-PR을 C-PR과 비교하면, 허락 화살표가 퍼뜨리는 길로서 보태는 것과 씨앗으로서 보태는 것을 나누어 볼 수 있어요.

결합 연산자에 드는 비용은 화살표 모음이 정해요. C-PR은 $E_A\cup E_T$ 위를 걷기 때문에, 반복 한 번에 두 층에서 반복을 한 번씩 하는 만큼의 비용이 들어요. 그래서 C-PR의 실행 시간은 AWP와 같거나 더 길 것으로 예상해요. 따라서 주장~C1의 속도 결과는 엔도스랭크만의 것이에요. S-PR은 층마다 한 번씩 풀어야 해요. 두 연산자 모두 \ref{sec:computational-benchmark}절의 벤치마크 절차에 따라 시간을 재요.

\dissertationSubheading[sec:holdout-protocol]{시간 홀드아웃 절차}

봄 홀드아웃(나중 기간을 떼어 두고 맞혀 보는 시험)에서 점수는 자료가 시작되는 2025년 12월 1일부터 $t_1=$ 2026년 2월 28일 23:59:59 UTC까지 꺼낸 모든 이벤트(계약이 남기는 기록)를 써요. 라벨(나중에 채점에 쓰는 정답 값)은 2026년 3월 1일 00:00:00 UTC부터 $t_2=$ 2026년 5월 31일 23:59:59 UTC까지의 이벤트로 세어요. 코호트(묶음)는 스펜더(spender, 허락을 받아 대신 쓰는 쪽) $n=1{,}335$개로 이루어져요. 이 스펜더들은 꺼낸 허락 자료에서 $t_1$에 0보다 큰 마지막 허락(가장 최근 허락)을 적어도 하나 가지고 있어요. 홀드아웃 표에 있는 점수는 모두 $t_1$까지의 이벤트만으로 계산해요. 그 점수들은 엔도스랭크(EndorseRank), AWP(적응형 가중 페이지랭크), 세 가지 $\lambda$ 값과 두 끝점에서의 C-PR(결합 페이지랭크, 두 그래프를 함께 걷는 점수), S-PR(보증 씨앗 페이지랭크), 그리고 원시 차수(화살표를 그냥 센 수) 기준선(비교 기준) 두 개예요. 시간 감쇠(오래된 일일수록 가볍게 세는 무게)의 기준점도 $t_1$이에요. 식~\eqref{eq:logistic-decay}의 $\Delta t$는 $t_1$에서 거꾸로 재요. 그래서 나이(이벤트가 일어난 뒤 지난 날수)는 $90$일보다 작고, 무게는 약 $0.71$과 $0.86$ 사이에 있어요.

지갑마다 라벨을 두 개 세어요. 구성개념(점수가 재려고 하는 생각)마다 하나씩이에요. 이렇게 하면 미래 기간에서, 한 층만 쓰는 방법은 저마다 자기 구성개념으로 시험받고, 혼합 점수(두 층을 섞은 점수)는 두 구성개념 모두로 시험받아요.
\begin{enumerate}
\item 새 허락 쌍(새로 생긴 주인–스펜더 허락): 라벨 기간에 0보다 큰 허락을 받았지만, $t_1$에 찍어 둔 ‘마지막 허락이 0보다 큰 쌍’의 기록(스냅숏, 그 순간의 모습)에는 없던 주인--스펜더 쌍의 수예요. 이 수는 스펜더 쪽에 매겨요. 이 라벨은 보증(믿고 맡김) 구성개념을 재요.
\item 새 송금자(새로 송금을 보내온 주소): 라벨 기간에 그 지갑에 송금(transfer)을 보낸 서로 다른 주소의 수예요. 다만 자료 시작일(2025년 12월 1일)과 $t_1$ 사이에는 그 지갑에 송금을 보낸 적이 없는 주소만 세고, 자기 자신에게 보낸 송금은 빼요. 이 라벨은 흐름(돈이 오가는 것) 구성개념을 재고, 첫 번째 라벨의 송금 쪽 짝이에요.
\end{enumerate}
기준선으로는 지갑마다 받은 허락 수(허락 입차수, 0보다 큰 허락을 해 준 서로 다른 주인의 수)와 송금 입차수(자기 자신을 뺀 서로 다른 송금자의 수)를 써요. 둘 다 $t_1$ 때의 값이에요. 한 층만 걷는 페이지랭크(PageRank)는 저마다 이 국소 통계(이웃만 보고 센 값) 가운데 하나를 매끄럽게 다듬어요(평활해요). 표에 이 기준선들의 줄도 있으므로, 각 페이지랭크가 자기 원시 차수보다 얼마나 더 나은지(증가분)를 바로 읽을 수 있어요. 연구 질문 RQ4에 답하는 것은 이 증가분이에요. 페이지랭크 계수 하나만으로는 답이 되지 않아요.

불확실성(결과가 얼마나 흔들릴 수 있는지)은 \ref{sec:statistical-tests}절의 짝지은 지갑 부트스트랩(지갑을 다시 뽑아 보기)으로 구해요($B=400$, 모든 점수가 함께 쓰는 다시 뽑기 번호 행렬 하나). 그래서 같은 라벨에서 두 계수 사이의 차이는 모두 짝지은 차이예요. 결합 점수를 계산하기 전에 정해 둔 대비(두 점수를 견준 차이)는 결합 연산자(두 층을 함께 걷게 하는 계산 규칙)와 그 층들에 관한 것이에요. 첫째는 새 허락 쌍에서 C-PR ($\lambda=1$) 빼기 받은 허락 수예요. 둘째는 새 허락 쌍에서 C-PR 빼기 C-PR ($\lambda=1$), 그리고 이와 함께 새 송금자에서 C-PR 빼기 C-PR ($\lambda=0$)이에요. 셋째는 S-PR에 대한 같은 두 대비예요. 엔도스랭크와 AWP가 들어가는 대비는 라벨을 안 뒤에 계산했어요. 각 점수를 자기 차수와 견준 대비와, C-PR을 그 둘 각각과 견준 대비예요. 엔도스랭크와 AWP 사이의 차이는 부록~\ref{ch:appendix-er-awp}에 있어요.

연구 질문 RQ5를 위한 첫 번째 규칙은 2026년 9월 14일에 정했어요. 이 규칙은 C-PR을, 자기 층을 하나만 따로 걷는 것과 견주어 판정하고, 두 부분으로 되어 있어요. 주된 부분은 홀드아웃을 써요. $\lambda=0.5$인 C-PR은 새 허락 쌍에서 C-PR ($\lambda=1$)보다 나쁘지 않아야 해요($\Delta\tau=\tau_{\mathrm{CPR}}-\tau_{\lambda=1}$의 구간(믿을 만한 범위)이 0을 포함하거나 0보다 위에 있어야 해요). 그리고 새 송금자에서는 C-PR ($\lambda=0$)보다 나아야 해요($\tau_{\mathrm{CPR}}-\tau_{\lambda=0}$의 구간이 0보다 위에 있어야 해요). 보조 부분은 같은 기간(점수를 계산한 바로 그 기간)과 매칭 코호트(모든 점수가 똑같이 점수를 매기는 지갑 묶음)를 써요. 그 조건은 두 가지예요. 송금 가족과 허락 가족(가족은 비슷한 지표 묶음)에서는 C-PR의 점 추정값(하나로 콕 집은 추정값)이 두 층 가운데 더 나은 쪽의 95\% 구간 안에 들어야 해요. 그리고 시빌 안정성(가짜 주소를 많이 만드는 공격에 흔들리지 않는 정도) 가족에서는 C-PR이 두 층을 모두 넘어야 해요. 표~\ref{tab:tau-diff-hybrid}의 짝지은 대비는 더 엄격한 점검으로 이 규칙 옆에 함께 보고해요. 하지만 그 대비가 규칙의 판정을 정하지는 않아요. 주된 부분이 실패하면, RQ5의 뒤쪽 절반, 곧 C-PR이 어느 한 층만 걷는 것보다 나은지에 대한 답은 “아니요”가 돼요. 그리고 그 결과를 구간과 함께 보고해요. 앞쪽 절반, 곧 C-PR이 송금만 걷는 것보다 나아지는지는 이 규칙이 다루지 않아요. 그것을 어떻게 시험하는지는 \ref{sub:fresh-holdout}절에서 설명해요. 라벨을 한 번 본 뒤에는 $\lambda$를 조정하지 않아요. 이 규칙을 정했을 때는 봄 라벨과, 한 층만 걷는 두 층의 홀드아웃 결과를 이미 알고 있었어요. 아직 계산하지 않은 것은 결합 점수뿐이었어요.

처리 과정(파이프라인)은 라벨 기간에 대해 세 가지를 더 계산했어요. 허락 취소 비율, 허락을 받은 뒤 토큰을 밖으로 빼내는지 보는 어림 규칙(빼내기 휴리스틱), 그리고 Aave(돈을 빌려주고 빌리는 블록체인 서비스)에서 돈을 빌린 사람의 청산(손실이 너무 커져 거래가 강제로 끝나는 일)이에요. 이 가운데 어느 것도 뒷받침하는 증거로 내놓지 않아요. \ref{ch:results}장에 허락 취소 계수와 빼내기 계수가 나와요. 이 계수들은 스펜더가 가진 주인 수가 많을수록 커져요. 그래서 믿음과 사용을 구별할 수 없어요. 이 스펜더 묶음에서 Aave 이벤트는 아예 비어 있거나, 엉뚱한 구성개념을 쟀어요.

\dissertationSubsubheading[sub:trader-holdout]{봄 기간의 트레이더 결과}

스펜더 코호트는 매칭 코호트를 이루는 트레이더(거래하는 사람)에 대해서는 알려 주는 것이 적어요. 트레이더 대부분은 허락을 받지 않기 때문이에요. 결합 연산자가 아직 없던 2026년 8월 29일에, 탐색적(미리 정하지 않고 살펴본) 실행 하나가 봄 동결일(점수를 고정하는 날)을 매칭 코호트에 적용했어요. 그리고 3월부터 5월까지 센 GMX(무기한 선물 거래 서비스) 성공 라벨을 골랐어요. 이익을 내며 닫은(포지션을 끝낸) 횟수, 실현 이익(실제로 손에 쥔 이익), 그리고 닫은 것 가운데 이익을 낸 것의 비율이에요. 코호트는 동결일 그래프(동결일 기준으로 만든 그래프)에 있는 매칭 지갑 4,227개예요. 라벨은 그 가운데 그 기간에 적어도 세 번 닫은 1,860개에 대해 정해져요. 엔도스랭크, AWP, C-PR은 라벨을 안 뒤에 이 코호트에서 점수를 매겼어요. 이때 스펜더 홀드아웃과 같은 짝지은 다시 뽑기 400번을 썼어요. \ref{ch:results}장은 이 실행을 탐색적인 것으로 보고해요. 이 실행에는 GMX 청산 라벨이 없어요. 그리고 Aave 청산 라벨은 다루는 지갑이 너무 적어서 해석할 수 없어요.

\dissertationSubsubheading[sub:fresh-holdout]{등록 재현: 2026년 6월부터 8월까지}

봄 홀드아웃은 어떤 규칙도 묻지 않았던 것을 하나 더 보여 주었어요. C-PR을 자기 송금 층만 걷는 것과 견주면, C-PR은 나중에 새 허락을 얻은 스펜더들을 훨씬 더 잘 줄 세웠어요. 그리고 새 송금자에서는 그 걷기와 맞먹었어요. 우리는 라벨이 다 들어온 뒤에야 이것을 알아챘어요. 그래서 이것은 발견이 아니라 가설로 쳐요. 그리고 이 시험을 정할 때 아직 기록을 꺼내지 않았던 기간에서 다시 시험해요. 설계는 2026년 9월 28일에 설정 파일 하나와 글로 쓴 계획 하나에 적었어요. 둘 다 기록을 꺼내기 전에 공개 저장소에 커밋(바뀐 내용을 기록으로 올리기)했어요(부록~\ref{ch:appendix-eval}). 2026년 10월 1일에는 같은 저장소에 커밋을 하나 올려, 설정 안의 상태를 ‘등록됨’(계획을 먼저 공개해 묶어 둠)으로 바꿨어요. 그리고 그날 늦게 6월부터 8월까지의 로그를 꺼냈어요. 꺼내기 스크립트는 두 파일이 커밋되고 상태가 ‘등록됨’으로 바뀌기 전에는 실행을 거부해요. 시험용으로 이 거부를 건너뛰는 장치가 있기는 해요. 그 장치를 쓰면 언제든 꺼내기 기록부(매니페스트)에 적히는데, 기록부에는 그런 기록이 하나도 없어요.

방법은 바뀌지 않았어요. C-PR은 9월 14일에 정한 대로 $\lambda=0.5$에서 쓰고, 두 민감도 값(결과가 얼마나 움직이는지 보려고 함께 쓰는 값)과 S-PR도 함께 써요. 풀이 프로그램(solver) 설정, SQL, 지갑 5,521개로 거르는 꺼내기 조건, GMX 해독기(기록을 읽어 내는 프로그램), 짝지은 부트스트랩(다시 뽑기 400번, 시드(무작위 뽑기를 똑같이 다시 할 수 있게 정해 둔 수) 42)은 봄에 쓴 것과 같아요. 점수 기간은 더 길어요. 점수는 2026년 5월 31일 23:59:59~UTC에 고정하고, 2025년 12월 1일부터 꺼낸 모든 이벤트를 써요. 그래서 송금과 허락의 나이가 90일이 아니라 182일까지 늘어나요. 라벨은 2026년 6월 1일부터 8월 31일까지 세어요.

두 코호트를 평가해요. 스펜더 코호트는 앞에서와 똑같이 정해요. 동결일에 0보다 큰 마지막 허락을 적어도 하나 가진 모든 스펜더예요. 트레이더 코호트는 매칭 코호트 가운데 동결일 그래프에 있는 지갑만 남긴 것이에요. 누가 여기에 드는지는 동결일 전에 한 닫기로 정해져요. 스펜더는 두 라벨을 그대로 가져가요. 트레이더는 청산 없이 닫은 비율을 라벨로 받아요. 이 비율은 1에서, 한 지갑이 6월부터 8월까지 닫은 것 가운데 청산으로 끝난 것의 비율을 뺀 값이에요. 그 기간에 적어도 세 번 닫은 지갑에 대해서만 정해요. 이 비율은 이 자료 안에서 채무 불이행(빌린 돈을 갚지 못하는 일)에 가장 가까운 것이에요. 또 어느 그래프로도 만들지 않았어요. 그리고 비율이기 때문에, 단지 거래를 더 많이 했다고 해서 지갑에 상을 주지 않아요. 이 라벨을 고를 때 이미 알고 있던 것이 두 가지 있어요. 하나는 청산이 얼마나 흔한지예요. 3월부터 5월까지, 적어도 세 번 닫은 매칭 지갑 2,962개 가운데 1,434개가 청산을 적어도 한 번 겪었어요. 그 가운데 동결일 그래프에 있던 1,860개 중에서는 923개가 그랬어요. 트레이더 코호트는 바로 이렇게 정해요. 다른 하나는 같은 기간에서의 짝이에요. 12월부터 5월까지 센 청산 없는 비율은 표~\ref{tab:alignment-summary}에 보관해 둔 청산 가족을 이루는 대리 지표(대신 재는 값) 세 개 가운데 하나예요. 이 지표 하나만 보면, 한 층만 걷는 허락 층은 $\tau=0.075$, 한 층만 걷는 송금 층은 $0.083$이 나왔어요. 이 비율은 C-PR에 대해서는 한 번도 계산한 적이 없었고, 기간 밖에서 계산한 적도 없었어요.

규칙은 C-PR을 자기 송금 층만 걷는 것, 곧 C-PR ($\lambda=0$)과 견줘요. 견주는 값은 짝지은 차이 $\Delta\tau=\tau_{\mathrm{CPR}}-\tau_{\lambda=0}$이에요.
\begin{itemize}
\item F1, 스펜더, 새 허락 쌍: $\Delta\tau$의 95\% 구간이 0보다 위에 있어야 해요.
\item F2, 스펜더, 새 송금자: 구간의 아래 끝이 $-0.02$보다 위에 있어야 해요.
\item F3, 트레이더, 청산 없이 닫은 비율: 구간의 아래 끝이 $-0.02$보다 위에 있어야 해요.
\end{itemize}
세 조건이 모두 맞을 때에만 C-PR이 송금만 걷는 것보다 나아진다고 쳐요. 그래서 여러 번 검사한 것을 바로잡는 보정(다중 검정 보정)은 필요 없어요. 허용 폭(이만큼까지는 뒤처져도 괜찮다고 정한 폭) 0.02는 봄 기간 F2 대비의 짝지은 구간의 절반 폭이에요. 이 표본 크기에서는 그보다 작은 차이를 가려낼 수 없었어요. F3의 약점 하나를 규칙과 함께 적어 두었어요. 두 점수 모두 청산을 미리 맞히지 못하면, 두 계수가 다 0 근처여도 F3은 통과할 수 있어요. 그래서 두 계수를 판정 옆에 함께 보고해요. 그 밖의 결과 세 개는 정해 둔 순서로 보고하고, 결정에는 들어가지 않아요. 첫째(F4)는 청산 없이 닫은 비율에서 C-PR을 C-PR ($\lambda=0$)과 견준 우월성 검사(더 나은지 보는 검사)예요. 둘째(F5)는 같은 라벨에서 C-PR을 C-PR ($\lambda=1$)과 견준 것이에요. 셋째(F6)는 9월 14일 규칙의 주된 부분을 새 기간에서 다시 해 본 것이에요. 등록 문서에서는 한 층만 걷는 두 층을 엔도스랭크와 AWP라고 불러요.

민감도 분석(조건을 바꿔 결과가 얼마나 움직이는지 보기) 두 개를 2026년 10월 1일에 더했어요. 봄 자료에서 AWP 점수를 매긴 뒤였고, 6월부터 8월까지의 로그를 하나라도 꺼내기 전이었어요. 둘 다 결정 옆에 보고하고, 어느 것도 결정에 들어가지 않아요. 첫 번째는 C-PR ($\lambda=0$) 대신 AWP(\ref{sub:awp-graph}절)를 넣고 F1--F4를 다시 해요. 두 번째는 그래프가 빠뜨린 코호트 지갑을 0점으로 매기는 대신 외톨이 점(화살표가 하나도 없는 점)으로 남겨 두고 F1--F6을 다시 해요(\ref{sec:rank-divergence}절). 두 번째 규칙에서 바뀌는 것은 한 층만 걷는 두 층뿐이에요. 모든 스펜더는 허락 그래프에 있고, 코호트의 모든 트레이더는 C-PR이 걷는 합친 그래프에 있기 때문이에요. 엔도스랭크는 등록한 점수에 들지 않았어요. 엔도스랭크는 등록한 평가가 끝난 뒤에 같은 코호트, 라벨, 지갑 다시 뽑기로 점수를 매겼어요. 그 결과는 결정과 따로 보고해요.

\dissertationSubheading[sec:proxy-evaluation]{구성개념 점검 설계}

같은 기간 평가에는 모든 점수가 들어가요. 엔도스랭크, AWP, 다섯 가지 $\lambda$ 값에서의 C-PR, 그리고 S-PR이에요. 처리 과정은 점수 $M$과 점검 $q$마다 $\mathbf{s}_M$과 $\mathbf{p}_q$ 사이의 스피어먼의 로($\rho$, 순위가 얼마나 함께 오르내리는지 나타내는 수)와 켄달의 타우($\tau$, 두 줄 세우기가 얼마나 같은 순서인지 나타내는 수)를 95\% 부트스트랩 구간과 함께 보고해요(\ref{sec:statistical-tests}절). AWP 검증 절차\cend{3}에서처럼, 점검마다 값이 클수록 기대되는 평판이 높다는 뜻이 되도록 방향을 맞춰요.

같은 기간 평가는 같은 시기의 가족 두 개를 중심으로 짜여 있어요.
\begin{enumerate}
\item 송금(AWP를 검증하는 데 쓴 대리 지표\cend{3}): 지갑으로 들어오는 송금 활동이에요.
\item 허락: 지갑이 스펜더일 때, 그 지갑으로 들어오는 허락 활동이에요.
\end{enumerate}
두 가족은 모두 그래프와 가까워요. 그 재료가 $G_{\text{transfer}}$와 $G_{\text{approve}}$를 이루는 이벤트 종류이기 때문이에요. 다만 점검 하나하나는 국소 통계이지, 그래프 전체를 보고 매긴 점수가 아니에요. AWP가 송금 점검과 아주 잘 맞거나 엔도스랭크가 허락 점검과 아주 잘 맞으면, 그것은 한 구성개념 안에서 앞뒤가 맞는다는 것을 보여 줄 뿐이에요. 따로 미리 맞히는 힘이 있다는 증거는 아니에요. 세 번째 가족은 \ref{sub:proxy-sybil}절의 시빌 보정 활동 요약이에요. 이 가족은 C-PR의 보조 규칙에 들어가지만, 이 가족에 기대는 가설은 없어요.

기간 밖 점검은 \ref{sec:holdout-protocol}절의 시간 홀드아웃이 맡아요. Packin과 Lev-Aretz\cend{2}는 탈중앙 신용 점수가 속을 들여다볼 수 없게(불투명하게) 될 수 있다고 경고해요. 계산식을 분명하게 공개하고(\ref{sec:proxy-formulas}절) 같은 기간 점검을 시간으로 나눈 시험과 따로 두면, 그런 불투명함이 줄어들어요.

\dissertationSubheading[sec:proxy-formulas]{대리 지표의 계산 정의}

아래에서 $\mathcal{W}$는 매칭 코호트를 뜻해요. 기호는 공개된 평가 처리 과정 안의 점검 계산 모듈(프로그램의 한 부분)에서 쓰는 것과 같아요.

\dissertationSubsubheading[sub:proxy-transfer]{송금 가족}

기간 안에 지갑 $w$가 받은 송금의 모음을 $\mathcal{T}_{\mathrm{in}}(w)$라고 쓰고, 송금 $e$를 보낸 쪽을 $u_e$라고 써요. 자기 자신에게 보낸 송금도 $\mathcal{T}_{\mathrm{in}}(w)$에 들어가요. 그래서 자기에게 토큰을 보낸 지갑은 자기 자신을 보낸 쪽 하나로 세어요. 양 $x_e$는 토큰을 가리지 않고 원래 단위(소수점 조정을 하지 않은 가장 작은 단위)로 더해요. 다음과 같이 정해요.
\begin{align}
\mathrm{in\_degree}(w)
&=
\bigl|\{\,u_e : e\in\mathcal{T}_{\mathrm{in}}(w)\,\}\bigr|,
\label{eq:in-degree}\\
\mathrm{in\_value}(w)
&=
\sum_{e\in\mathcal{T}_{\mathrm{in}}(w)} x_e.
\label{eq:in-value}
\end{align}
둘 다 AWP 논문\cend{3}에서 정한 기준선이에요. 경제적으로 중요한 주소일수록 더 많은 서로 다른 거래 상대에게서, 더 큰 총액의 송금을 받는다는 생각을 바탕으로 해요. 이 둘은 페이지랭크와 달리 국소적이고, 보낸 쪽의 평판은 따지지 않아요. 그래서 AWP를 이 둘과 견주면, 퍼뜨림(전파), 최근 일에 무게를 더 주기, 활발하게 보내는 주소에서 다시 출발하기가 AWP의 순서를 그냥 들어온 양의 순서에서 얼마나 멀리 옮기는지 알 수 있어요.

\dissertationSubsubheading[sub:proxy-allowance]{허락 가족}

허락 쪽에서는 $w$가 스펜더예요(주인 $\to w$). 식~\eqref{eq:approve-weight}의 마지막 허락을 $a^{\star}$라고 하고, 다음과 같이 정해요.
\begin{align}
\mathrm{in\_approve\_degree}(w)
&=
\bigl|\{\,u : a^{\star}(u,w,\cdot)>0\,\}\bigr|,
\label{eq:in-approve-degree}\\
\mathrm{in\_approve\_value}(w)
&=
\sum_{u,\mathrm{token}} a^{\star}(u,w,\mathrm{token}).
\label{eq:in-approve-value}
\end{align}
이 점검들은 엔도스랭크가 순위를 매기도록 만들어진 바로 그 보증을 재요. 그래서 둘의 순서가 아주 잘 맞는 것은 예상된 일이고, 구성개념 안의 타당도(같은 구성개념 안에서 서로 맞는지)로 읽어요. 송금 기준선처럼, 이 점검들도 들어오는 허락을 세고 더할 뿐, 여러 단계를 거치는 보증 사슬을 따라가지는 않아요.

\dissertationSubsubheading[sub:proxy-sybil]{시빌 보정 안정성 가족}

들어오는 것을 세는 값에서는 자기 자신에게 보낸 송금($u=w$)을 빼요. $w$가 다른 주소에게서 받은 송금의 수를 $N_{\mathrm{tx}}(w)$라 하고, 들어오는 쪽의 서로 다른 거래 상대 수를 $N_{\mathrm{uniq}}(w)$라 해요. 머문 기간(첫 기록과 마지막 기록 사이의 날수)과 활동한 달 수는 두 번째 줄 묶음을 써요. 코호트 지갑이 낀 송금은 모두, 보낸 쪽이 코호트에 있으면 보낸 쪽의 것으로 치고, 그렇지 않으면 받는 쪽의 것으로 쳐요. 그래서 코호트 지갑 두 개 사이의 송금은 보낸 쪽에게만 세어져요. $w$의 것으로 친 줄들 가운데 가장 이른 시각과 가장 늦은 시각(타임스탬프)을 $t_{\min}(w)$와 $t_{\max}(w)$로 잡아요. $N_{\mathrm{mon}}(w)$는 서로 다른 달력 달을 센 두 수 가운데 더 큰 것이에요. 하나는 $w$가 다른 주소에게서 송금을 받은 달의 수이고, 다른 하나는 $w$의 것으로 친 줄들이 있는 달의 수예요. 다음과 같이 정해요.
\begin{align}
\mathrm{inbound\_counterparty\_ratio}(w)
&=
\frac{N_{\mathrm{uniq}}(w)}{\max(N_{\mathrm{tx}}(w),1)},
\label{eq:counterparty-ratio}\\
\mathrm{transfer\_tenure\_days}(w)
&=
\frac{t_{\max}(w)-t_{\min}(w)}{86400}\quad\text{(초를 날로 바꿈)},
\label{eq:tenure}\\
\mathrm{active\_months}(w)
&=
N_{\mathrm{mon}}(w).
\label{eq:active-months}
\end{align}
거래 상대 비율은 적은 주소에게서 송금을 많이 받는 지갑에 벌점을 줘요. 이것은 단순한 자전 거래(자기끼리 사고팔아 거래가 많은 척하기) 무늬예요. 하지만 흔한 시빌 무늬, 곧 새로 만든 주소 여러 개가 송금을 하나씩 보내는 무늬에는 벌점을 주지 않아요. 그렇게 송금을 받은 지갑은 가장 큰 값인 1에 가까운 점수를 받아요. 머문 기간과 활동한 달 수는 관찰 기간 동안 온체인(블록체인 위의) 활동이 얼마나 꾸준한지를 따라가요. 그러므로 이 대리 지표들은 약한 시빌 점검이에요. 서로 짠 공격자들은 세 가지를 모두 만족시킬 수 있어요. 그래도 이 지표들은, 오래 이어지고 여러 상대와 하는 활동이 한꺼번에 몰아서 하는 자기 거래보다 평판에 대해 더 많은 것을 말해 준다는 직관을 계산할 수 있는 모양으로 옮겨 놓아요\cend{2}.

\dissertationSubheading[sec:statistical-tests]{통계 검사}

평판 점수도, 검증용 대리 지표도 채무 불이행인지 아닌지 둘 중 하나만 나타내는 값이 아니에요. 둘 다 연속적인 값이고, 순위로 다뤄요. AWP 검증 절차\cend{3}와 구성 타당도(점수가 재려는 것을 정말 재는지)에 대한 Cronbach와 Meehl의 접근 방식\cend{22}을 따라, 순서가 맞는 정도는 순위 상관(순서가 얼마나 함께 움직이는지)으로 재요. 순위 상관은 두 변수 가운데 어느 쪽이든 순서가 바뀌지 않게 값을 다시 늘리거나 줄여도 바뀌지 않아요. 계수마다 부트스트랩 신뢰구간(믿을 만한 범위)을 함께 보고해요.

\begin{itemize}
\item 스피어먼의 $\rho$: 지갑 $i$의 원래 점수와 원래 대리 지표 값의 중간 순위 $R_i$와 $S_i$로 구한 피어슨 상관(두 값이 곧은 줄을 따라 함께 움직이는 정도)이에요. 값이 같으면 그 자리 번호들의 평균을 함께 나눠 가져요(\ref{sec:rank-correlation-theory}절). 이 수는 단조 일치(한 방향으로 함께 움직이는지)를 재요. 평판이 높을수록 대리 지표 값도 늘 높으면, 원래 점수가 곧은 줄에서 얼마나 벗어나든 $\rho$는 양수예요.

\item 켄달의 $\tau$: 보고하는 값은 모두 식~\eqref{eq:kendall}의 동점(점수가 같은 것)을 바로잡은 $\tau_b$이고, 원래 점수와 대리 지표로 계산해요. 동점이 없으면 $(1+\tau)/2$는 아무렇게나 고른 한 쌍을 두 순위가 같은 순서로 놓을 확률이에요. 동점이 있으면 이 풀이는 $\tau_b$가 아니라 $\tau_a=(n_c-n_d)/n_0$에 맞는 것이에요. 그런데 여기서는 동점이 흔해요. 매칭 코호트에서 모든 지갑 쌍의 $67.3\%$가 엔도스랭크에서 동점이에요.

\item 둘 다 쓰는 까닭: AWP 절차\cend{3}가 둘 다 보고해요. 그리고 둘 다 보고하면 결론이 계수 하나에만 매달리지 않아요. 점수 분포는 꼬리가 두꺼워요(아주 큰 값이 드물지 않아요). 대부분의 지갑은 0 가까이에 있고, 허브(화살표가 많이 모이는 점) 몇 개가 점수의 큰 몫을 차지해요. 그래서 원래 점수로 구한 피어슨 상관 대신 순위 계수를 써요. 피어슨 상관은 허브들이 좌우해 버리기 때문이에요.

\item 지표 가족 합치기: 지표 가족 하나는 그 안의 대리 지표들을 단순 평균해서 요약해요. 예를 들어 송금 가족의 평균 $\tau$는 $\tfrac{1}{2}(\tau_{\mathrm{in\_degree}}+\tau_{\mathrm{in\_value}})$예요. 정해진 코호트와 정해진 대리 지표에서 점수끼리 견준 것은 그 점수들에 대한 상대적인 말일 뿐이에요. 평판의 보편적인(어디서나 통하는) 성질을 뜻하지는 않아요. 우리의 결론은 신뢰구간에서 끌어내요. 어떤 절대적인 $\tau$ 기준값도 결정 기준으로 쓰지 않아요.

\item 불확실성: 표본을 뽑는 데서 오는 불확실성은 짝지은 지갑 부트스트랩\cend{34}으로 어림해요. 매칭 코호트를, 뽑은 것을 다시 넣어 가며 $B=400$번 다시 뽑아요(시드 42, 설정에 고정). 그리고 번호 행렬 하나를 모든 방법과 대리 지표에 다시 써요. 그래서 다시 뽑은 통계들은 방법끼리 짝지어져 있어요. 대리 지표마다의 $\tau$, 지표 가족 평균, 방법 사이 차이마다, 부트스트랩 분포의 2.5번째와 97.5번째 백분위수(줄 세웠을 때 그 퍼센트 자리에 오는 값)가 95\% 백분위 구간을 이뤄요\cend{35}. $\tau_b$는 U-통계량(쌍들을 견주어 평균한 값)의 매끄러운 함수이므로, $n=5{,}521$에서 백분위 구간은 1차까지 정확해요. 다시 뽑는 횟수가 적어서, 동점을 바로잡은 $\tau$의 $O(B\,n\log n)$ 계산 비용을 모든 방법과 대리 지표에 걸쳐 감당할 수 있어요. 그 대신 구간의 끝값은 약 $\pm 0.005$까지만 정확해요.

\item 차이 검사: 모든 비교는 같은 다시 뽑은 표본 위에서 하는 짝지은 대비예요. 다시 뽑은 표본마다 $\Delta\tau=\tau_A-\tau_B$를 계산하고, 그 부트스트랩 분포의 95\% 백분위 구간을 보고해요. 구간이 0을 포함하지 않으면 그 대비를 차이로 쳐요. 다시 뽑기가 짝지어져 있으므로, 대비는 두 추정값의 공분산(두 값이 함께 움직이는 정도)을 그대로 지녀요. $\operatorname{Var}(\hat\tau_A-\hat\tau_B)=\operatorname{Var}\hat\tau_A+\operatorname{Var}\hat\tau_B-2\operatorname{Cov}(\hat\tau_A,\hat\tau_B)$예요. 그래서 공분산이 양수이면, 대비의 구간은 따로 구한 두 구간을 보고 짐작하는 것보다 좁아요. 규칙의 판정을 정하는 대비들은, 판정에 쓰는 자료보다 먼저 정해 두었어요. 9월 14일 규칙의 대비(\ref{sec:holdout-protocol}절)와 등록 재현(계획을 먼저 공개해 두고 새 자료로 다시 해 보기)의 대비(\ref{sub:fresh-holdout}절)예요. 봄 기간 새 허락 쌍에서 C-PR을 자기 송금 층만 걷는 것과 견준 대비는 여기에 들지 않았어요. 그래서 이 대비는 나올 때마다 탐색적이라고 표시해요. 같은 기간에서 대비 여덟 개는 C-PR과 S-PR을, 가족마다 앞서는 층(한 층만 걷는 것)과 견줘요. 송금 가족에서는 송금 걷기, 허락 가족에서는 허락 걷기, 시빌 안정성에서는 둘 다예요. 대비 네 개는 C-PR을 송금에서는 AWP와, 허락에서는 엔도스랭크와, 시빌 안정성에서는 둘 다와 견줘요. 엔도스랭크와 AWP가 들어가는 같은 기간 대비 다섯 개는 어느 것도 미리 정하지 않았고, 부록~\ref{ch:appendix-er-awp}에 보고해요.

\item 방향 맞추기와 빈 값 채우기: 대리 지표마다 값이 클수록 평판이 좋다는 뜻이 되도록 방향을 맞춰요. $\mathcal{W}$에 다시 맞춘 뒤, 어떤 대리 지표의 자료에 없는 지갑은 0으로 채워요. 거래가 듬성듬성하므로 이것은 조심스러운(보수적인) 선택이에요. 그리고 두 점수 매기기 방식 모두 이 방법을 똑같이 써요. 이 단계를 거치면 코호트의 모든 지갑이 모든 점수 벡터와 대리 지표 벡터에 값을 가져요. 그래서 상관마다 $\mathcal{W}$ 전체를 써요.

\item 비열등성(정한 폭보다 더 나쁘지 않음): 등록 재현의 세 조건 가운데 두 개는, C-PR이 자기 송금 층만 걷는 것보다 미리 정한 허용 폭 $\delta=0.02$보다 더 많이 뒤처지지 않기만을 요구해요. 이런 대비는 95\% 구간의 아래 끝이 $-\delta$보다 위에 있으면 통과해요. 이 검사는 차이 검사보다 약하고, 그 라벨에서 더 나은지(우월성)에 대해서는 아무것도 말해 주지 않아요. 이 검사는 C-PR이 중요한 것을 아무것도 내주지 않는다는 주장을 시험할 때 써요.

\item 빼놓은 것: 이 연구는 원인과 결과를 밝히는 추론(인과 추론)을 하지 않아요. 따로 나눈 학습용 자료와 시험용 자료로 지도 학습 모델(정답을 보고 배우는 모델)을 훈련하지도 않고, $p$값도 보고하지 않아요. 대신 미리 정한 대비의 구간이 ‘차이 없음’에 해당하는 값을 포함하지 않는지를 물어요.

\end{itemize}

\dissertationSubheading[sec:computational-benchmark]{계산 벤치마크(주장~C1)}

주장~C1은 점수마다 드는 계산 비용에 관한 것이에요. 벤치마크(속도 재기)는 화살표(엣지) 목록을 점수로 바꾸는 풀이 프로그램 호출에 걸리는 시간을 재요. 화살표 목록을 만드는 데 한 번만 드는 비용은 빼요. 이 비용은 대부분 불러오기와 전처리(미리 손보기)이고, 점수를 정기적으로 새로 고치는 점수 서비스라면 여러 번에 나누어 부담하게 될 비용이에요.

\begin{itemize}
\item 시간을 재는 부분: 화살표 목록과 AWP의 재시작 무게는 시간 재기를 시작하기 전에 만들어요. 시간을 재는 호출은 점(노드)에 번호를 매기고, 듬성듬성한 무게를 단 인접 행렬(어느 점이 어느 점으로 이어지는지 적은 표)을 조립하고, 줄마다 합이 1이 되게 나누고, 재시작 벡터를 만들고, $d=0.85$, $\varepsilon=10^{-8}$, $I_{\max}=300$으로 거듭제곱 반복(같은 계산을 되풀이하기)을 돌려요. 같은 호출을 따로 분석해 본 실행(어디에 시간이 드는지 재 보기)에서는 행렬 조립이 실제 걸린 시간의 약 94\%, 거듭제곱 반복이 약 1\%를 차지했어요. 두 단계 모두 $O(|E|)$이고, 실제 서비스에 쓴다면 갱신할 때마다 두 단계를 모두 다시 해야 해요. 일치도 점수도 같은 그래프에서 나와요.

\item 측정 절차: 모든 설정은 한 세션 안에서 재요. 먼저 시간을 재지 않는 준비 실행을 한 번 하는데, 가장 많이 쓴 메모리는 이 실행에서만 기록해요. 그다음 시간을 재는 실행을 다섯 번 해요. 이 다섯 번에서 실행 시간의 평균과 표준편차(SD), 평균 반복(한 바퀴) 횟수를 입력 그래프의 점과 화살표 수와 함께 보고해요. 결과는 화살표 다중집합(같은 것이 여러 번 들어갈 수 있는 모음), 감쇠 계수(화살표를 따라갈 확률), 허용 오차, 반복 상한(가장 많이 반복할 수 있는 횟수), 몇 번째 실행인지로 만든 열쇠(키) 아래 저장해 두고 다시 써요. 같은 설정의 같은 그래프가 여러 곳에 나와요. 주 벤치마크, 감쇠 계수 바꿔 보기의 $d=0.85$ 줄, $n=5{,}521$인 코호트 안 규모 단계, 그리고 전체 풀 단계예요. 전체 풀 단계의 유도 부분 그래프(고른 점들과 그 사이 화살표만 남긴 그래프)가 바로 코호트 그래프예요(\ref{sub:expanded-pool}절). 그래서 그 표들은 일부러 같은 값을 보고하며, 한 양을 따로따로 잰 독립된 측정이 아니에요. 반복 한 번에 $O(|E|)$만큼 일하는 알고리즘에서는 화살표가 비용을 좌우하므로\cend{15}, 모든 실행 시간 표에 화살표 수를 적어요. 그리고 실행 시간은 컴퓨터가 읽을 수 있는 출력에 적어요(\ref{sec:reproducibility}절).

\item 코호트: 주 벤치마크는 매칭 지갑 코호트($n=5{,}521$)에서 돌려요. 점 개수에 대한 민감도를 보려고, 확장 지갑 풀(더 큰 주소 모음)에서 늘 똑같이 나오는 작은 표본(결정적 부분 표본)을 뽑아요(\ref{sec:scaling-robustness}절).

\item 공정성: 모든 알고리즘은 같은 풀이 프로그램, 허용 오차, 반복 상한, 감쇠 계수(감쇠 계수 바꿔 보기는 빼고), 씨앗 지갑 묶음(그래프를 만들 때 출발점이 되는 지갑들)으로 돌아가요. 그래서 비용의 차이는 구현에서 오지 않아요. 그래프에서, 그리고 반복 횟수를 바꿀 수 있는 AWP의 재시작 벡터에서 와요. 하드웨어와 소프트웨어는 평가 기록부(매니페스트)에 적어요. 두 점수의 비용은 시간 비율로 견주고, 시간 비율마다 그에 맞는 화살표 수 비율을 옆에 적어요.

\end{itemize}

\dissertationSubheading[sec:scaling-robustness]{규모에 따른 변화와 강건성으로 넓힌 분석}

우리는 더 큰 주소 모음에서 실행 시간이 점 개수에 따라 어떻게 바뀌는지도 살펴봐요. 또 지표 가족 수준의 무늬가 감쇠 계수, 토큰 모음, 표본을 바꿔도 버티는지도 살펴봐요. 이 분석들은 주장~C1에 단서를 달고, 지표 가족 수준 결론이 어디까지 통하는지 시험해요. 어느 쪽도 주 표본을 대신할 다른 표본을 내놓지 않아요. 따로 말하지 않으면, \ref{ch:results}장의 일치도에 관한 주장은 매칭 코호트에 대한 것이고, 점 개수에 대한 민감도는 확장 풀에 대한 것이에요.

\dissertationSubsubheading[sub:expanded-pool]{확장 지갑 풀과 점 개수에 따른 실행 시간의 민감도}

확장 지갑 풀 $\mathcal{W}_{\mathrm{pool}}$은 매칭 코호트(주소 $5{,}521$개)와 추가 주소 $93{,}559$개를 합친 것으로 정해요. 추가 주소는 관찰 기간 동안 달마다의 아비트럼 원(Arbitrum One) ERC-20(토큰 공통 규칙) 활동에서 뽑았어요. 우리 표본 뽑기 스크립트가 만드는 추가 활동 지갑 parquet 파일(표 모양 자료를 담는 파일 형식)에서 가져온 것이에요. 그래서 이 연구에는 $N=|\mathcal{W}_{\mathrm{pool}}|=99{,}080$개의 주소가 있어요. 중간 크기 벤치마크는 시드 문자열 \texttt{benchmark-tier2-v1}을 쓴 결정적 SHA256 작은 표본 뽑기로 풀에서 뽑아요. 지갑들은 다음 값으로 줄을 세워요.
\begin{center}
\texttt{SHA256(}\textit{시드}\texttt{:}\textit{주소}\texttt{)}, 여기서 \textit{시드} $=$ \texttt{benchmark-tier2-v1},
\end{center}
곧, 시드와 쌍점(:)과 소문자로 바꾼 주소를 이어 붙인 문자열 하나의 해시(글자를 뒤섞어 만든 고정 길이 값)로 줄을 세워요. 그리고 그 순서에서 처음 $n$개 지갑을 작은 표본으로 삼아요. 시드 문자열은 우리 설정에 고정해 둔 상수예요. 세 단계를 보고해요. $n=10{,}000$, $n=50{,}000$, 그리고 전체 풀 $n=N=99{,}080$이에요. 모두 같은 허락·송금 parquet 파일을 써요. $n\le 5{,}521$인 코호트 안 실행은 부록~\ref{ch:appendix-eval}에 있어요. 해시 계산은 결정적이므로, 의사 난수(컴퓨터가 흉내 낸 무작위 수) 뽑기와 달리 라이브러리 버전이 바뀌어도 작은 표본이 달라질 수 없어요.

이 설계의 두 가지 특징 때문에, 이 실험이 보여 줄 수 있는 것에는 한계가 있어요. 첫째, 모든 단계에서 그래프는 코호트 기간의 허락·송금 parquet 파일에 있는 화살표 가운데, 끝점이 적어도 하나 뽑힌 지갑 안에 있는 것만 남겨요. 그리고 풀이 프로그램이 쓰는 점은 그 화살표들의 끝점이에요. 그래서 어떤 단계의 화살표 수 $|E|$는 해시 순서가 코호트 지갑을 몇 개 고르는지에 달려 있어요. 곧 이 수는 $n$을 따라가지 않아요. 두 방법 모두에서, $n=10{,}000$인 풀 단계는 $n=2{,}000$인 코호트 안 단계보다 화살표가 적어요. 다른 특징은 추가 풀 주소에 관한 것이에요. 이 주소들 자신의 허락·송금 로그는 두 번째 단계에서 꺼낼 계획이었어요. 그런데 그 꺼내기가 끝나지 않았어요(부록~\ref{ch:appendix-eval}). 그 결과, 전체 풀 단계는 코호트 parquet 파일 밖의 화살표를 하나도 더하지 않고, 코호트 그래프를 그대로 다시 만들어요. 점도 하나도 더하지 않아요. 남긴 화살표에 하나도 닿지 않는 뽑힌 주소는 풀이 프로그램의 그래프에 아예 들어가지 않기 때문이에요. 전체 풀의 허락 그래프에는 코호트의 점 $6{,}442$개가 그대로 있어요. 그러므로 풀 단계들은 외톨이 점이 더해질 때 실행 시간이 어떻게 바뀌는지 재지 않아요. 또 반복 한 번의 비용을 좌우하는 화살표 수를 키우지도 않아요. 규모에 따른 변화를 보여 주는 표는 모두 두 방법에 대해 단계마다 $|V|$와 $|E|$를 적어요. 그리고 \ref{ch:results}장에서 규모에 따른 변화에 관해 하는 말은 그 열들이 허락하는 데까지만 나아가요.

\dissertationSubsubheading[sub:damping-protocol]{감쇠 계수 민감도}

매칭 코호트에서 그래프와 대리 지표는 그대로 두고, 감쇠 계수를 $d\in\{0.75,0.85,0.95\}$로 바꿔 봐요. $d$가 작으면 순간 이동(다시 출발)이 많아져요. 반대로 크면 여러 단계를 거치는 길에 무게를 더 주고, 수렴(값이 더는 거의 바뀌지 않게 됨)이 느려질 수 있어요. 이렇게 바꿔 보기는 허락 가족과 송금 가족에서의 일치도가 순간 이동 확률에 달려 있는지를 물어요.

\dissertationSubsubheading[sub:top-tokens]{상위 토큰 부분 그래프}

이 점검에서는 ERC-20 토큰 $20$개만 남기고, 그 부분 그래프에서 엔도스랭크와 AWP를 다시 계산해요. 첫 번째 고르기는 기간 동안 합친 양이 가장 큰 토큰들을 골라요. 이때 원래 송금 양과 원래 허락 값을 더해요. 양은 원래 단위이고 무제한 허락은 $2^{256}-1$로 세어요. 그래서 이 순위는 무제한 허락이 많고 공급량(만들어진 토큰의 양)이 큰 토큰에 유리해요. 결과를 안 뒤에 더한 두 번째 고르기는 줄(송금 이벤트와 마지막 허락을 합친 것)이 가장 많은 토큰들을 골라요. 어느 쪽도 경제적 가치로 매긴 순위는 아니에요. 이 점검은 나머지 토큰을 모두 빼도 지표 가족 수준의 결과가 바뀌는지를 물어요.

\dissertationSubsubheading[sub:sample-size-protocol]{표본 정의 민감도}

단계마다, 같은 결정적 시드로 뽑은 확장 풀의 작은 표본에서 일치도 통계를 다시 계산해요. 모든 일치도 주장의 주 표본은 그대로 매칭 코호트이고, 그 계수와 구간도 매칭 코호트에서 어림해요. 풀의 작은 표본은 크기와 구성이 코호트와 달라요. 그 안의 지갑 가운데 많은 수가 한 그래프에만 나와요. 그래서 다른 그래프에서의 점수와 대리 지표가 0이에요. 그리고 0에서 생기는 동점은 순위 일치를 저절로 부풀려요. 그러므로 단계마다의 계수는 표본 정의에 대한 민감도 점검이지, 안정성의 증거가 아니에요. 묻는 것은 표본이 바뀔 때 계수가 저마다 얼마나 움직이는지예요.

\dissertationSubheading[sec:reproducibility]{재현성}

\ref{ch:results}장의 표는 하나도 손으로 옮겨 적지 않았어요. 1단계(Phase~1)에서 꺼낸 자료가 디스크에 있으면, 평가 실행 프로그램(드라이버) 하나가 모든 점검, 구간이 붙은 모든 일치도 통계, 모든 벤치마크를 다시 계산해요. 그리고 그것들을 컴퓨터가 읽을 수 있는 요약에 적어요. 이 요약에서 따로 있는 단계가 \ref{ch:results}장과 부록~\ref{ch:appendix-eval}의 표를 내보내요. 이 실행 프로그램은 저장해 둔 추출 자료나 합성 시험 자료(만들어 낸 연습용 자료)를 입력으로 받아요. 그리고 감쇠 계수, 상위 토큰, 표본 정의를 바꿔 보는 스위치가 있어요. 버전을 매긴 설정 파일 하나에 관찰 기간, 감쇠 계수, 시간 감쇠의 매개변수, 이벤트 토픽 ID, 부트스트랩 횟수와 시드, 벤치마크 시드가 들어 있어요. 시험 자료 모드에서는 처리 과정이 클라우드 접속 자격 증명 없이, 인터넷에 연결하지 않고도 돌아가요.

다시 해 볼 수 있는 요소는 이래요. SHA256 해시로 하는 결정적인 지갑 작은 표본 뽑기, 고정된 페이지랭크 허용 오차와 반복 상한, 부트스트랩 시드 하나와 다시 뽑기 번호 행렬 하나, \ref{sec:proxy-formulas}절의 대리 지표 식, 실행마다 잰 시간을 로그에 남기는 \ref{sec:computational-benchmark}절의 벤치마크 절차, 실행을 시작하기 전에 둘 다 설정 파일에 고정한 섞는 비율 $\lambda$와 \ref{sec:holdout-protocol}절의 홀드아웃 논리, 그리고 꺼낸 자료의 기록부(매니페스트) 파일이에요. 바로 그 추출 자료에서, 새 빅쿼리(BigQuery) 질의 없이 일치도 표와 벤치마크 표를 다시 만들 수 있어요. 모든 표 뒤에 있는 컴퓨터가 읽을 수 있는 요약은, 그때 쓴 설정과 함께 Git 저장소의 \texttt{data/3-published-results-for-thesis/} 아래에 커밋되어 있어요. 소스 코드와 같은 폴더예요. 그래서 처리 과정을 다시 돌리지 않고도 이 논문의 숫자를 그 요약과 맞춰 볼 수 있어요. 부록~\ref{ch:appendix-eval}에는 GitHub 저장소 주소(URL), 이 논문의 표를 만든 정확한 커밋 해시, 설정 파일, 요약 파일이 적혀 있어요.

\dissertationSubheading[sec:sampling-bias]{표본 뽑기에서 생각할 점과 빠진 자료}

매칭 코호트는 기준에 따라 만든 표본이지, 모든 아비트럼 주소에서 무작위로 뽑은 것이 아니에요. 코호트에 들려면 기간 안에 GMX~V2에서 적어도 세 번 닫아야 하고, 허락 활동과 송금 활동은 여기에 아무 역할도 하지 않아요. 그러므로 같은 기간 결과는 아비트럼 원에서 활발한 GMX~V2 트레이더에게 해당하고, 홀드아웃 결과는 이 트레이더들의 허락 자료에 있는 스펜더에게 해당해요. 기간 안에 GMX~V2에서 거래하지 않은 지갑은 설계상 빠져요. 잠자는 지갑(오래 쓰지 않은 지갑)과 토큰을 오래 들고만 있는 사람도 여기에 들어요. 그리고 그래프 밖에 있는 코호트 지갑은 그 방법에서 $0$점을 받아요(\ref{sec:data-sources}절).

자료는 몇 가지 방식으로 빠질 수 있어요. 읽어 낼 수 없는 양은 0으로 읽고, 값이 0인 기록과 함께 버려요. 엔도스랭크는 표준 ERC-20 이벤트를 남기지 않는 허락을 보지 못해요. AWP는 오프체인(블록체인 밖)에서, 또는 보관 계정(거래소 같은 곳이 대신 맡아 관리하는 계정) 안에서 한 송금을 보지 못해요. 그래서 둘 다 온체인, 표준 토큰 행동 쪽으로 치우쳐 있어요. 이것은 온체인 평판 연구로서는 받아들일 만하지만, 결과를 한 서비스의 모든 사용자에게 넓혀 적용할 때는 문제가 돼요.

풀($N=99{,}080$)은 실행 시간 실험에서 뽑는 주소 모음을 넓혀 줘요. 하지만 풀이 프로그램의 그래프는 코호트 그래프보다 커지지 않아요(\ref{sub:expanded-pool}절). 그래서 일치도는 언제나 매칭 코호트에서 계산해요.

\dissertationSubheading[sec:implementation-notes]{구현 메모}

페이지랭크 코드는 Python으로 썼어요. SciPy로 만든 듬성듬성한 이동 연산자와 NumPy로 하는 거듭제곱 반복을 써요. 설정마다 화살표 목록은 한 번만 만들어요. 시간 재기 벤치마크는 입출력(I/O)과 전처리를 빼고, 연산자 조립과 반복만 다뤄요(\ref{sec:computational-benchmark}절). 무작위성이 들어오는 곳은 골라서 쓰는 시험 자료 만들기 한 군데뿐이에요. parquet 입력과 설정의 시드가 주어지면, 진짜 자료에 대한 평가는 결정적이에요. 플랫폼(컴퓨터 환경)마다 부동소수점(컴퓨터가 소수를 다루는 방식) 계산에서 작은 차이가 생기는 것은 예상할 수 있는 일이에요. 하지만 그 차이는 소수 셋째 자리까지 보는 $\tau$ 지표의 정밀도에 비하면 무시할 만큼 작을 거예요.

그림은 표와 같은 평가 결과물로 만들어요. 그래서 그림에 그린 값은 표에 적힌 숫자와 맞아요.

\dissertationSubheading[sec:ethics]{윤리}

이 연구는 누구나 볼 수 있는 블록체인 자료만 쓰고, 사람을 연구 대상으로 삼지 않아요. 기관의 윤리 심사(Unisa SoC/CAES)는 실제 자료를 모으기 시작하기 전에 끝났어요. 부록~\ref{sec:ethics-certificate}에 윤리 승인 증명서가 있어요(참조 번호 [Unisa SoC/CAES 윤리 승인 참조 번호: TBD], 날짜 [TBD]).

\begin{itemize}
\item\phantomsection\label{sub:privacy-anonymity}
  사생활과 익명성: 아비트럼 원의 주소는 가명이에요. 주소만으로는 그 뒤에 있는 실제 사람이 누구인지 드러나지 않아요\cend{2}. 하지만 가명성(진짜 이름 대신 주소로만 활동함)은 익명성(누구인지 전혀 알 수 없음)과 같지 않아요. 이론상으로는 마음먹은 공격자가 거래소 입금, 온라인 프로필, 어림 규칙을 이용해 주소를 한 사람과 이어 붙일 수 있어요. 이 연구에서는 일부러 그런 연결을 하나도 하지 않아요. 우리의 자료 모으기는 온체인 \texttt{Approval} 이벤트와 \texttt{Transfer} 이벤트로만 한정돼요. 그리고 오프체인의, 개인을 알아볼 수 있는 정보(IP 주소, 거래소 KYC(고객 신원 확인) 기록, 소셜 미디어 아이디, 익명을 벗기는 그래프 같은 것)는 하나도 긁어 오지 않았어요. 우리는 주소들을 실제 세상의 주체로 묶지 않았어요. 개인을 알아볼 수 있을 것 같은 주소의 순위표를 공개하지도 않았어요. 신원을 알아내려고 지갑의 거래 내역을 다른 데이터베이스와 합쳐 보지도 않았어요. 공개한 결과는 모아서 계산한 값(상관 계수, 실행 시간, 지표 가족 평균)뿐이에요. 이름이 밝혀진 참가자 한 명 한 명의 순위는 들어 있지 않아요. 본문에서 예를 들 때는 넓은 사용자 분류(예를 들어 ‘허락 금액이 큰 스펜더’)만 말하고, 특정 주소는 절대 말하지 않아요.


\item\phantomsection\label{sub:data-minimization}
  자료 최소화(꼭 필요한 자료만 갖기): 처리한 표에는 그래프 만들기와 점검에 필요한 칸만 들어 있어요. 곧 관련 주소, 토큰 양, 시각 기록(타임스탬프)이에요. 원래 로그는 연구 처리 과정을 감사(잘못이 없는지 검사)할 수 있도록 보관해요. 사용자의 익명을 벗기려는 목적으로는 공개하지 않을 거예요. 가능한 곳에서는 질의의 거름 조건을 지갑으로 좁혔어요. 연구 설계의 범위를 벗어나, 체인 전체를 가리지 않고 통째로 쏟아 내는 일을 피하려는 것이에요.


\item\phantomsection\label{sub:data-usage-compliance}
  자료 사용과 규칙 지키기: 모든 질의는 \texttt{bigquery-public-data.goog\_blockchain\_arbitrum\_one\_us} 자료 모음에 올라 있는 공개 자료에 대해 돌렸어요. 그리고 구글 클라우드(Google Cloud)의 이용 약관과 사용량 제한을 지켰어요. 미리 돌려 보기로 얻은 바이트 추정값과 질의마다의 최대치 덕분에 필요 없는 훑기를 막았어요. 우리는 거래를 하나도 보내지 않았고, 특별 권한이 있거나 공개되지 않은 색인 접속 지점에도 접속하지 않았어요. 또 네트워크의 정상 작동을 방해할 수 있는 방식으로 스마트 계약(블록체인 위에서 저절로 실행되는 프로그램)을 부른 일도 전혀 없어요.


\item\phantomsection\label{sub:fairness-transparency}
  공정성과 투명성: 새로 온 참가자, 가끔만 활동하는 사람, 공개된 허락을 하나도 해 주지 않는 지갑은 모두 그래프로 평판을 매기는 시스템에서 낮게 평가될 수 있어요\cend{2}. 엔도스랭크와 AWP 둘 다 이런 구조적인 치우침을 가지고 있어요. 예를 들어 허락 그래프에서는, 들어오는 쓰기 권한을 하나도 받지 못한 주소가 받은 주소보다 불리해요. 한편 송금 그래프는 돈이 흐르는 그물망 안에 박혀 있는 주소를 편들어요. \ref{ch:discussion}장에서 한계를 다뤄요. 그 가운데에는 처음 시작하는 주소가 불리한 점(콜드 스타트 효과)과, 허락 그래프가 시빌 공격(가짜 주소를 많이 만들어 속이는 공격)에 약한 점이 있어요. 모든 코드 경로, 매개변수 선택, 평가 지표는 공개된 처리 과정에 기록되어 있어요(부록~\ref{ch:appendix-eval}). 그래서 다른 사람이 이 연구를 검토하고 따로 다시 해 볼 수 있어요.


\item\phantomsection\label{sub:responsible-disclosure}
  책임 있는 공개와 영향: 엔도스랭크는 디파이(DeFi, 탈중앙 금융)의 위험 평가에 참고가 될 수 있어요. 하지만 금융 조언도, 신용 조언도, 소비자 신용 보고 상품도 아니에요. 돈을 빌려주는 일, 거래, 규칙 준수에 관한 결정에는 여전히 사람의 감독이 필요해요. \ref{ch:discussion}장에서는 공격자가 조작할 가능성을 인정하고 다뤄요. 예를 들어 $G_{\text{approve}}$에 대한 시빌 공격이나 $G_{\text{transfer}}$에 대한 자전 송금이에요. 이 연구는 코호트 수준에서 모아서 계산한 결과를 보고하고, 개인 사용자를 겨냥하려는 지갑 수준의 순위는 보고하지 않아요. 앞으로 실제로 서비스에 쓴다면 운영 규칙(거버넌스), 이의 제기 절차, 공정성 감시가 더 필요할 거예요. 이것들은 모두 이 연구의 범위 밖이에요.


\item\phantomsection\label{sub:methodological-commitments}
  방법에 대한 약속 요약: 결과를 해석하기 전에, 설계는 이미 체인, 기간, 코호트 규칙, 그래프 식, 대리 지표 정의, 부트스트랩과 시간 재기 절차, 통계 검사를 정해 두었어요. 주된 주장은 모든 점수를 함께 다뤄요. 매칭 지갑 코호트와 스펜더 홀드아웃에서, 각 페이지랭크를 그것이 매끄럽게 다듬는 원시 차수와 견줘요. 확장 풀과 민감도 바꿔 보기는 효율, 그리고 결과가 매개변수와 표본 정의에 얼마나 달려 있는지를 시험해요.
\end{itemize}

\ref{ch:results}장은 이 절차가 아비트럼 원에서 내놓는 결과를 보고해요. 자료와 코호트, 실행 시간 벤치마크, 모든 점수의 같은 기간 일치도, 순위 갈라짐, 민감도 바꿔 보기, 시간 홀드아웃들과 등록 재현, 그리고 마지막으로 주장~C1--C5예요.
